\chapter{Transformer——注意力就是一切}
\label{ch:transformer}
%% ============================================================

2017 年，Google 的研究团队发表了一篇改变历史进程的论文：
《Attention is All You Need》~\cite{vaswani2017attention}。
这篇论文提出的 Transformer 架构，在不到十年的时间里，
从一个机器翻译模型发展成了当今所有 AI 系统的核心引擎。
为什么 Transformer 如此成功？让我们从最核心的组件——注意力机制——开始理解。

\section{自注意力：学会「看哪里」}

\subsection{直觉：图书馆的类比}

想象你走进一座图书馆，你有一个问题想要回答。
你的问题就是\textbf{查询}（Query）：
「我想找关于法国历史的信息。」
图书馆里每本书的标题是\textbf{键}（Key）：
「法国大革命」「量子物理导论」「巴黎美食指南」……
每本书的内容是\textbf{值}（Value）。

你会把自己的查询与每本书的标题做比较，
与「法国大革命」的匹配度最高，
与「巴黎美食指南」也有一定相关性，
与「量子物理导论」则几乎不相关。
然后你根据这些相关性，加权地阅读各本书的内容。

这个类比可以更精确地对应到信息检索领域。
在搜索引擎中，用户输入查询 $q$，
系统计算 $q$ 与每个文档键 $k_i$ 的相关性分数 $\mathrm{score}(q, k_i)$，
然后按分数排序返回文档内容 $v_i$。
注意力机制做的事情完全一样，
只不过它不是返回 top-$k$ 文档，
而是对\textbf{所有}文档做\textbf{软性}加权平均：
相关性越高的文档权重越大。

自注意力（Self-Attention）做的就是这件事。
在一个句子中，\textbf{每个词同时扮演查询、键和值三个角色}。
每个词都在问：「句子里的其他词中，哪些与我最相关？」
然后根据相关性，加权地整合其他词的信息。

\subsection{更深入的直觉：软性字典查找}

图书馆类比可以被精炼为一个更精确的概念：
\textbf{注意力是一种可微分的软性字典查找}
（differentiable soft dictionary lookup）。

传统字典查找是硬性的：给定键 $k$，
要么找到完全匹配的条目并返回对应的值，
要么返回空。这是一个离散的、不可微的操作。
注意力机制将这个过程变成了连续的、可微的：
查询 $q$ 不是精确匹配某一个键，
而是\textbf{同时与所有键计算相似度}，
然后按相似度加权返回所有值的混合。

这个视角揭示了自注意力的本质能力：
\textbf{内容寻址的动态信息路由}。
每个 token 根据自身的\emph{内容}（而非固定的位置或规则）
决定从序列中的哪些位置提取信息。
这与传统编程中的 hash map 查找截然不同：
hash map 要求精确匹配，而注意力允许模糊匹配、
部分匹配、甚至在语义空间中的跨概念匹配。

\subsection{Q/K/V：同一信息的三种视角}

Q、K、V 三个投影矩阵的存在常常令初学者困惑：
为什么需要三个不同的投影？
不能直接用原始向量计算相似度吗？

答案在于\textbf{角色分离}。
考虑一个具体的例子：句子「The cat sat on the mat because it was tired.」
对于代词「it」，模型需要解决指代——「it」指的是「cat」还是「mat」？

如果直接用原始嵌入计算相似度，
「it」与「cat」和「mat」的相似度取决于词嵌入空间中的距离：
这主要反映的是词义相似度（两个名词可能距离相近），
而非指代关系。

Q/K/V 投影的作用是\textbf{将同一信息投影到不同的功能子空间}：
\begin{itemize}[noitemsep]
  \item $W_Q$ 将 token 投影到「我在找什么信息」的子空间：
        「it」作为查询时，其投影向量编码了
        「我需要找到一个可以作为主语的名词性实体」
  \item $W_K$ 将 token 投影到「我能提供什么信息」的子空间：
        「cat」作为键时，其投影向量编码了
        「我是一个有生命的名词性实体」
  \item $W_V$ 将 token 投影到「我实际携带的信息」的子空间：
        「cat」作为值时，其投影向量编码了
        关于猫的语义特征
\end{itemize}

这种三重投影允许模型学习到：
在查询--键空间中，「it」与「cat」高度匹配
（因为代词指代的是有生命的主语），
同时在值空间中传递「cat」的完整语义信息。
如果不做投影分离，这种灵活的匹配--传递机制就无法实现。

更形式化地说，Q/K 投影的组合 $W_Q^\top W_K$
定义了一个\textbf{双线性匹配函数}
$\mathrm{score}(x_i, x_j) = x_i^\top W_Q^\top W_K x_j$，
它可以学到输入空间中\emph{任意线性可分的相关性模式}。
V 投影则独立控制信息传递的内容。
三者各司其职。

\subsection{高维空间中的点积：几何直觉}

点积在高维空间中有一些反直觉的性质，
理解这些性质对于理解注意力至关重要。

在 $d$ 维空间中，两个向量 $\mathbf{q}$ 和 $\mathbf{k}$ 的点积：
\[
  \mathbf{q} \cdot \mathbf{k} = \|\mathbf{q}\| \|\mathbf{k}\| \cos\theta
\]
其中 $\theta$ 是两者之间的夹角。
如果两个向量是独立随机生成的，
点积的期望值为 0，方差为 $d$
（每个维度贡献约 1 的方差）。
这意味着\textbf{点积的绝对值随维度增长而增大}：
当 $d = 1024$ 时，随机点积的标准差约为 $\sqrt{1024} = 32$。

这就是为什么\textbf{缩放因子 $\sqrt{d_k}$} 如此关键。
不缩放的点积在高维中会产生很大的绝对值，
而 softmax 对大数值极其敏感：
当输入值远离零时，softmax 的输出趋近于 one-hot 向量，
梯度趋近于零，模型陷入「过度自信」的状态而无法学习。
除以 $\sqrt{d_k}$ 将点积的方差归一化为 1，
使 softmax 保持在梯度有效的区间内。

另一个重要的几何现象是\textbf{高维空间中的近似正交性}。
当 $d$ 足够大时，任意两个随机向量几乎一定近似正交
（$\cos\theta \approx 0$）。
这意味着在未经训练的初始状态下，
所有 token 对之间的注意力分数几乎相等
（因为点积都接近 0）：
这是一个良好的起始点，
让模型从「均匀关注」出发逐渐学习出有区分度的注意力模式。

\subsection{数学公式}

给定输入序列中某个位置的向量 $\mathbf{x}$，
我们用三个可学习的线性变换将其映射为查询、键、值：
\begin{equation}
  \mathbf{q} = W_Q \mathbf{x}, \quad
  \mathbf{k} = W_K \mathbf{x}, \quad
  \mathbf{v} = W_V \mathbf{x}
\end{equation}
其中 $W_Q, W_K \in \mathbb{R}^{d_k \times d}$，
$W_V \in \mathbb{R}^{d_v \times d}$ 是参数矩阵。

对整个序列做矩阵运算，注意力的计算公式为：
\begin{equation}
  \mathrm{Attention}(Q, K, V)
  = \mathrm{softmax}\!\left(\frac{Q K^\top}{\sqrt{d_k}}\right) V
  \label{eq:attention}
\end{equation}

这个公式分为三步。
第一步，$QK^\top$ 计算所有查询与所有键的点积，
得到一个 $n \times n$ 的相关性矩阵（$n$ 是序列长度）。
第二步，除以 $\sqrt{d_k}$ 进行缩放：
这一步至关重要，因为当维度 $d_k$ 较大时，
点积的值会变得很大，
导致 softmax 的梯度趋近于零（饱和），模型无法有效学习。
第三步，softmax 将相关性分数转换为概率分布，
然后与 $V$ 相乘得到加权输出。

\subsection{计算复杂度分析}

注意力的计算复杂度值得仔细分析，
因为它决定了 Transformer 能处理多长的序列。

$QK^\top$ 的计算：$Q \in \mathbb{R}^{n \times d_k}$，
$K \in \mathbb{R}^{n \times d_k}$，
矩阵乘法的复杂度为 $O(n^2 d_k)$。
softmax 的复杂度为 $O(n^2)$（对 $n \times n$ 矩阵逐行操作）。
最终 $\mathrm{softmax}(\cdot) \times V$ 的复杂度为 $O(n^2 d_v)$。
总计算复杂度为 $O(n^2 d)$，其中 $d = d_k = d_v$。

内存方面，需要存储 $n \times n$ 的注意力矩阵。
对于 $n = 128{,}000$（128K 上下文），
这个矩阵有 $128K \times 128K \approx 160$ 亿个元素：
仅以 FP16 存储就需要约 30\,GB。
这就是为什么长上下文是 LLM 的核心挑战之一，
也是 Flash Attention 如此重要的原因。

\subsection{手工计算示例}

为了让注意力的计算过程完全透明，
让我们用一个极简的例子手动走一遍。
假设序列长度 $n = 3$，每个 token 的维度 $d = 4$，
查询/键/值的维度 $d_k = d_v = 2$。

输入矩阵（3 个 token，每个 4 维）：
\[
X = \begin{pmatrix}
  1 & 0 & 1 & 0 \\
  0 & 1 & 0 & 1 \\
  1 & 1 & 0 & 0
\end{pmatrix}
\]

投影矩阵（简化为固定值）：
\[
W_Q = \begin{pmatrix} 1 & 0 \\ 0 & 1 \\ 1 & 0 \\ 0 & 1 \end{pmatrix}, \quad
W_K = \begin{pmatrix} 0 & 1 \\ 1 & 0 \\ 0 & 1 \\ 1 & 0 \end{pmatrix}, \quad
W_V = \begin{pmatrix} 1 & 1 \\ 0 & 0 \\ 0 & 1 \\ 1 & 0 \end{pmatrix}
\]

\textbf{Step 1}: 计算 $Q, K, V$：
\[
Q = XW_Q = \begin{pmatrix} 2 & 0 \\ 0 & 2 \\ 1 & 1 \end{pmatrix}, \quad
K = XW_K = \begin{pmatrix} 0 & 2 \\ 2 & 0 \\ 1 & 1 \end{pmatrix}, \quad
V = XW_V = \begin{pmatrix} 1 & 2 \\ 1 & 0 \\ 1 & 1 \end{pmatrix}
\]

\textbf{Step 2}: 计算 $QK^\top / \sqrt{d_k}$（$\sqrt{d_k} = \sqrt{2} \approx 1.41$）：
\[
QK^\top = \begin{pmatrix} 0 & 4 & 2 \\ 4 & 0 & 2 \\ 2 & 2 & 2 \end{pmatrix}
\quad \Rightarrow \quad
\frac{QK^\top}{\sqrt{2}} \approx \begin{pmatrix} 0 & 2.83 & 1.41 \\ 2.83 & 0 & 1.41 \\ 1.41 & 1.41 & 1.41 \end{pmatrix}
\]

\textbf{Step 3}: 对每一行做 softmax，得到注意力权重矩阵 $A$：
\[
A \approx \begin{pmatrix}
  0.05 & 0.77 & 0.19 \\
  0.77 & 0.05 & 0.19 \\
  0.33 & 0.33 & 0.33
\end{pmatrix}
\]

\textbf{Step 4}: 加权求和 $\mathrm{Output} = AV$：
\[
\mathrm{Output} \approx \begin{pmatrix}
  1.00 & 0.28 \\
  1.00 & 1.72 \\
  1.00 & 1.00
\end{pmatrix}
\]

观察结果：token 1 的输出主要由 token 2 的值决定（权重 0.77），
token 2 的输出主要由 token 1 的值决定（权重 0.77），
而 token 3 均匀地关注所有 token（权重各约 0.33）。
这就是注意力的本质，让每个 token 选择性地聚焦于最相关的其他 token。

\subsection{代码实现}

用 PyTorch 实现自注意力仅需几行代码：

\begin{lstlisting}[caption={Self-Attention in PyTorch}]
import torch
import torch.nn.functional as F
import math

def self_attention(x, W_q, W_k, W_v):
    """
    x: (batch, seq_len, d_model)
    W_q, W_k, W_v: (d_model, d_k) or (d_model, d_v)
    """
    Q = x @ W_q  # (batch, seq_len, d_k)
    K = x @ W_k  # (batch, seq_len, d_k)
    V = x @ W_v  # (batch, seq_len, d_v)

    d_k = Q.shape[-1]
    # Step 1+2: scaled dot-product
    scores = (Q @ K.transpose(-2, -1)) / math.sqrt(d_k)
    # Step 3: softmax + weighted sum
    weights = F.softmax(scores, dim=-1)
    output = weights @ V
    return output
\end{lstlisting}

这段代码对应了公式~\eqref{eq:attention}的三个步骤。
\texttt{Q @ K.transpose(-2, -1)} 是 $QK^\top$（矩阵乘法），
除以 \texttt{math.sqrt(d\_k)} 是缩放，
\texttt{F.softmax} 是 softmax 归一化，
最后 \texttt{weights @ V} 是加权求和。

\subsection{注意力的矩阵视角：信息流的全局图景}

从矩阵运算的角度审视注意力，可以获得更宏观的理解。

注意力矩阵 $A = \mathrm{softmax}(QK^\top / \sqrt{d_k})$
是一个 $n \times n$ 的\textbf{行随机矩阵}
（每行的元素非负且求和为 1）。
它本质上定义了一个\textbf{加权有向图}：
$A_{ij}$ 表示 token $i$ 从 token $j$ 接收信息的权重。

这个图结构揭示了注意力的信息流模式：
\begin{itemize}[noitemsep]
  \item 如果 $A$ 接近单位矩阵，
        每个 token 只关注自己——信息不流动
  \item 如果 $A$ 的某些列有很大的值，
        对应的 token 是「信息枢纽」：
        被许多其他 token 关注（如 [CLS] token 或句号）
  \item 如果 $A$ 接近均匀分布（每个元素约 $1/n$），
        所有 token 的信息被等权平均：
        这相当于一个全局池化操作，丢失了位置信息
\end{itemize}

从线性代数角度看，输出 $\mathrm{Output} = AV$
是对 $V$ 矩阵的行做\textbf{凸组合}。
$V$ 的每一行是一个 token 的值向量，
输出的每一行是所有值向量的加权平均。
这意味着注意力输出\textbf{永远在输入值向量的凸包内}：
注意力不能创造新信息，只能\emph{重新组合}现有信息。
新信息的创造由残差连接和 FFN 层完成。

这个凸组合性质也解释了为什么\textbf{深层堆叠}是必要的。
单层注意力只能做一次凸组合：
能力有限。多层注意力的堆叠
允许信息在层之间逐步整合和精炼，
每一层在前一层的输出基础上做进一步的组合，
最终实现复杂的信息变换。
Sanford 等人~\cite{sanford2023representational}
在理论上证明了 $L$ 层 Transformer 能表达的函数类
严格大于 $L-1$ 层——每增加一层都真正扩展了表达能力。

\begin{whybox}[为什么用点积而不是其他相似度度量？]
Bahdanau 等人在 2015 年提出的注意力使用加法形式
$\mathrm{score} = v^\top \tanh(W_1 q + W_2 k)$，
需要额外的参数和非线性操作。
点积注意力 $q^\top k$ 没有额外参数，
而且可以直接用高度优化的矩阵乘法实现：
这在 GPU 上的效率远高于逐元素操作。
在实践中，缩放点积注意力的效果与加法注意力相当，
但速度快得多。
此外，点积在几何上有清晰的含义：
$q^\top k = \|q\| \|k\| \cos\theta$，
即两个向量的相似度由它们之间的夹角决定。
这使得注意力权重具有直觉上的可解释性：
内容越相似（夹角越小），注意力越高。
\end{whybox}

\section{多头注意力：同时关注不同关系}

单个注意力头只能学习一种关注模式。
但在语言中，词与词之间存在多种关系：
语法关系（主语--谓语）、语义关系（同义词、反义词）、
指代关系（代词指向哪个实体）、位置关系（相邻的词往往相关）。

多头注意力（Multi-Head Attention, MHA）的解决方案是：
\textbf{并行运行多个注意力头，每个头独立学习不同的关注模式}。

\begin{equation}
  \mathrm{MultiHead}(Q, K, V) =
  \mathrm{Concat}(\mathrm{head}_1, \ldots, \mathrm{head}_h) \, W^O
\end{equation}
其中 $\mathrm{head}_i = \mathrm{Attention}(Q W_i^Q, K W_i^K, V W_i^V)$。

每个头使用自己的投影矩阵（$W_i^Q, W_i^K, W_i^V$），
在不同的子空间中计算注意力。
最后将所有头的输出拼接起来，通过一个线性层 $W^O$ 融合。

\subsection{为什么多头比单头好？}

直觉上，多头注意力的好处是\textbf{分工}：
不同的头可以专注于不同类型的语言关系。
Clark 等人~\cite{clark2019bert}
在对 BERT 的注意力模式进行可视化分析后发现，
不同的头确实学到了截然不同的模式：

\begin{itemize}[noitemsep]
  \item 某些头专注于\textbf{局部上下文}（如相邻词），
        类似于 n-gram 模型
  \item 某些头关注\textbf{语法依赖}，
        如主语和谓语之间的一致性关系，
        即使它们相距很远
  \item 某些头追踪\textbf{指代关系}：
        「他」指向哪个人物，「它」指向哪个对象
  \item 某些头负责\textbf{分隔符注意力}：
        大量注意力集中在句号、逗号等标点符号上，
        可能用于分割句子结构
\end{itemize}

从数学角度看，多头注意力的总参数量与单头注意力相同。
如果模型维度为 $d$，头数为 $h$，
那么每个头的维度为 $d_k = d/h$。
单头注意力的投影矩阵大小为 $d \times d$，
多头注意力中每个头的投影矩阵大小为 $d \times d/h$，
$h$ 个头加起来仍然是 $d \times d$。
因此，多头注意力\textbf{不增加参数量}，
但通过在不同子空间中并行计算，获得了更丰富的表示能力。

现代 LLM（如 LLaMA 3、DeepSeek-V3）通常使用 32--128 个注意力头。

\subsection{多头注意力的可视化}

图~\ref{fig:mha-architecture}展示了多头注意力的完整数据流，
从输入到多个并行的注意力头再到最终的融合输出。

\begin{figure}[htbp]
\centering
\begin{tikzpicture}[
  node distance=10mm,
  box/.style={draw, rounded corners=2.5pt, minimum height=10mm,
    font=\small, inner sep=5pt, line width=0.7pt},
  io/.style={box, fill=figLightGray, draw=figGray!60, minimum width=60mm},
  proj/.style={box, fill=figBlue!8, draw=figBlue!50, minimum width=18mm},
  head/.style={box, fill=figOrange!8, draw=figOrange!50, minimum width=24mm},
  merge/.style={box, fill=figGreen!8, draw=figGreen!50},
  splitbar/.style={box, fill=figLightGray, draw=figGray!40,
    minimum width=68mm, minimum height=6mm, font=\scriptsize},
  arr/.style={-{Stealth[length=2.5pt,width=2.5pt]}, line width=0.6pt, figGray!60},
  lab/.style={font=\scriptsize, text=figGray},
]

% Input
\node[io] (input) {$X \in \mathbb{R}^{n \times d}$};

% Projections
\node[proj, above=12mm of input, xshift=-22mm] (wq) {$W^Q$};
\node[proj, above=12mm of input] (wk) {$W^K$};
\node[proj, above=12mm of input, xshift=22mm] (wv) {$W^V$};
\node[lab, above=1mm of wq] {Query};
\node[lab, above=1mm of wk] {Key};
\node[lab, above=1mm of wv] {Value};

% Input -> Projections
\draw[arr] (input.north -| wq) -- (wq);
\draw[arr] (input.north) -- (wk);
\draw[arr] (input.north -| wv) -- (wv);

% Split/Reshape bar (replaces the 9-arrow fan-out)
\node[splitbar, above=8mm of wk] (split)
  {Reshape: $(n,\, d) \to (n,\, h,\, d_k)$};

% Projections -> Split (3 clean vertical arrows)
\draw[arr] (wq.north) -- (wq.north |- split.south);
\draw[arr] (wk.north) -- (split.south);
\draw[arr] (wv.north) -- (wv.north |- split.south);

% Heads
\node[head, above=10mm of split, xshift=-24mm] (h1) {Head$_1$};
\node[head, above=10mm of split] (h2) {Head$_2$};
\node[head, above=10mm of split, xshift=24mm] (hh) {Head$_h$};
\node[font=\small, text=figGray] at ($(h2.east)!0.5!(hh.west)$) {$\cdots$};

% Head formula (side annotation)
\node[lab, right=6mm of hh, text width=36mm, align=left] {%
  $\textstyle\text{head}_i = \text{softmax}\!\bigl(\frac{Q_i K_i^\top}{\sqrt{d_k}}\bigr)\! V_i$\\[4pt]
  $d_k = d/h$};

% Split -> Heads (3 clean vertical arrows)
\draw[arr] (h1.south |- split.north) -- (h1.south);
\draw[arr] (split.north) -- (h2.south);
\draw[arr] (hh.south |- split.north) -- (hh.south);

% Concat
\node[merge, above=12mm of h2, minimum width=58mm] (concat)
  {Concat};
\draw[arr] (h1) -- (h1 |- concat.south);
\draw[arr] (h2) -- (concat);
\draw[arr] (hh) -- (hh |- concat.south);

% Output projection
\node[merge, above=8mm of concat, minimum width=28mm] (wo) {$W^O$};
\draw[arr] (concat) -- (wo);

% Output
\node[io, above=8mm of wo] (output) {Output $\in \mathbb{R}^{n \times d}$};
\draw[arr] (wo) -- (output);

\end{tikzpicture}
\caption{多头注意力数据流。输入 $X$ 分别经过 $W^Q$, $W^K$, $W^V$ 线性投影，
切分为 $h$ 个并行注意力头，各头在 $d/h$ 维子空间独立计算缩放点积注意力，
拼接后经 $W^O$ 投影为 $d$ 维输出。总参数量与单头注意力一致。}
\label{fig:mha-architecture}
\end{figure}

\subsection{注意力头的冗余与剪枝}

一个有趣的发现是：并非所有注意力头同等重要。
Michel 等人~\cite{michel2019heads}
在对 BERT 和 NMT 模型的分析中发现，
在许多层中，移除大部分注意力头后模型性能几乎不变：
某些层甚至只需要保留\textbf{一个}头。

这一发现暗示了两点：
第一，标准 MHA 存在显著的\textbf{参数冗余}，
多个头可能学到了高度重叠的模式。
第二，关键的注意力模式只存在于少数头中，
大部分头起到的是辅助或冗余的作用。

这种冗余性恰好为 MQA 和 GQA 的有效性提供了理论支撑：
如果许多头本来就在学习相似的 KV 模式，
让它们共享 KV 不会造成太大的信息损失。
同时，这也是结构化剪枝和知识蒸馏在注意力层特别有效的原因。

\subsection{从 MHA 到 MQA 到 GQA：效率的演进}

MHA 的问题在推理时暴露无遗。
在自回归生成中，每生成一个新 token，
都需要用到之前所有 token 的 $K$ 和 $V$。
为了避免重复计算，这些 $K, V$ 会被缓存：
这就是 \textbf{KV cache}。

对于标准 MHA，每个头都有独立的 $K$ 和 $V$ 矩阵。
一个 70B 参数模型处理 128K token 的上下文时，
KV cache 的大小为：
\[
  2 \times n_{\mathrm{layers}} \times n_{\mathrm{heads}}
  \times \mathrm{seq\_len} \times d_{\mathrm{head}} \times \mathrm{bytes}
\]
代入典型值（80 层、64 头、128K 序列、128 维、FP16），
$2 \times 80 \times 64 \times 128{,}000 \times 128 \times 2 \approx 335$\,GB：
远超单块 GPU 的显存。

\textbf{Multi-Query Attention（MQA）}~\cite{shazeer2019mqa}
是 Shazeer 在 2019 年提出的激进方案：
\textbf{所有查询头共享一套 $K, V$}。
KV cache 从 $h$ 套降到 1 套，内存减少到 $1/h$。
代价是质量下降：
所有头被迫使用相同的 KV 信息，表达能力受限。

\textbf{分组查询注意力}（Grouped-Query Attention, GQA）
是 Ainslie 等人（2023）提出的折中方案~\cite{ainslie2023gqa}，被 LLaMA 2 广泛采用~\cite{dubey2024llama3}：
将多个查询头分成若干组，每组共享一套 KV。
例如，32 个查询头分成 8 组，只需要 8 套 KV：
KV cache 减少到原来的 $1/4$，
而质量几乎没有损失。

GQA 处于两个极端之间：

\begin{center}\small
\renewcommand{\arraystretch}{1.3}
\begin{tabular}{lccc}
\toprule
\rowcolor{figLightGray}
\textbf{方法} & \textbf{KV 头数} & \textbf{KV cache 大小} & \textbf{质量} \\
\midrule
MHA  & $h$（全部独立） & $1\times$ & 最高 \\
GQA  & $h/g$（分组共享） & $1/g$ & 接近 MHA \\
MQA  & 1（全部共享） & $1/h$ & 可能下降 \\
\bottomrule
\end{tabular}
\end{center}

GQA 在质量和效率之间找到了最优平衡，
是目前几乎所有主流 LLM 的标准选择。
LLaMA 3 使用 $h = 32, g = 4$（即 8 组 KV 头），
DeepSeek-V2 则走得更远，采用了 MLA（下文详述）。

\subsection{滑动窗口注意力}

标准注意力让每个 token 关注序列中的所有 token，
但在实践中，大部分注意力权重集中在附近的 token 上。
\textbf{滑动窗口注意力}（Sliding Window Attention, SWA）
利用了这一观察，每个 token 只关注距自己最近的 $w$ 个 token，
其中 $w$ 是窗口大小。

SWA 的核心优势是将每个 token 的注意力计算
从 $O(n)$ 降为 $O(w)$，总复杂度从 $O(n^2)$ 降为 $O(nw)$。
当 $w \ll n$ 时，这是巨大的节省。

Mistral 7B~\cite{jiang2023mistral} 是 SWA 的标志性应用。
Mistral 使用 $w = 4096$ 的滑动窗口，
配合 32 层的堆叠：
虽然每层只能看到 4K 的局部上下文，
但信息可以通过残差连接在层之间传播。
经过 $L$ 层后，理论上的\textbf{有效感受野}
为 $L \times w$（$32 \times 4096 = 128K$ tokens）。

这种设计的精妙之处在于：
局部信息在底层通过 SWA 高效处理，
全局信息通过层间传播隐式地整合。
它在不改变注意力计算公式的前提下，
通过限制注意力掩码来实现加速。

SWA 可以与 GQA 组合使用：
Mistral 同时采用了两者。
SWA 减少了每个头的计算量，
GQA 减少了 KV cache 的内存。

\subsection{线性注意力：突破二次瓶颈}

标准注意力的 $O(n^2)$ 复杂度源于
$QK^\top$ 的矩阵乘法需要计算所有 $n^2$ 个元素对。
\textbf{线性注意力}~\cite{katharopoulos2020linear}
通过改变计算顺序来避免这个瓶颈。

标准注意力（忽略 softmax 中的归一化常数）可以写成：
\[
  \mathrm{Output}_i = \frac{\sum_j \exp(q_i^\top k_j / \sqrt{d_k}) \cdot v_j}
                           {\sum_j \exp(q_i^\top k_j / \sqrt{d_k})}
\]

线性注意力的关键是用\textbf{核函数}
$\phi(\cdot)$ 近似 softmax 中的指数函数：
\[
  \exp(q^\top k / \sqrt{d_k}) \approx \phi(q)^\top \phi(k)
\]
替换后，注意力变成：
\[
  \mathrm{Output}_i = \frac{\phi(q_i)^\top \sum_j \phi(k_j) v_j^\top}
                           {\phi(q_i)^\top \sum_j \phi(k_j)}
\]

关键在于 $\sum_j \phi(k_j) v_j^\top$ 和 $\sum_j \phi(k_j)$
可以\textbf{预先计算并缓存}（维度为 $d' \times d_v$ 和 $d'$），
不需要为每个查询重新遍历所有键值对。
这样，计算复杂度从 $O(n^2 d)$ 降为 $O(n d d')$，
其中 $d'$ 是核函数的输出维度，通常远小于 $n$。

但线性注意力有一个根本性的代价：
它无法精确表达标准 softmax 注意力的全部功能。
特别是，softmax 的\textbf{尖峰特性}
（能将大部分注意力集中在少数 token 上）
在核近似下会被平滑化：
模型失去了精确聚焦的能力。
这就是为什么纯线性注意力的模型在实践中
通常弱于标准注意力模型的原因。

线性注意力的思想深刻影响了后续的 SSM 和递归架构设计。
实际上，Mamba-2 的 SSD 理论~\cite{dao2024mamba2}
正是建立在线性注意力与 SSM 之间的数学对偶性上。
从这个意义上说，线性注意力不是失败的尝试，
而是通向高效序列建模的理论桥梁。

\subsection{注意力变体的统一视角}

回顾以上各种注意力变体，可以从三个正交维度理解它们的设计空间：

\textbf{维度一：KV 共享策略}（减少内存）
\begin{itemize}[noitemsep]
  \item MHA：每个头独立 KV → 无共享
  \item GQA：每组头共享 KV → 组内共享
  \item MQA：所有头共享 KV → 全局共享
  \item MLA：通过低秩投影压缩 KV → 压缩共享
\end{itemize}

\textbf{维度二：注意力范围}（减少计算）
\begin{itemize}[noitemsep]
  \item 全局注意力：$O(n^2)$，精确但昂贵
  \item 滑动窗口：$O(nw)$，局部高效，依赖层间传播
  \item 稀疏注意力（NSA）：$O(n \cdot s)$，动态选择 $s$ 个关键 token
  \item 线性注意力：$O(nd)$，通过核近似消除二次项
\end{itemize}

\textbf{维度三：注意力质量}（提升精度）
\begin{itemize}[noitemsep]
  \item 标准 softmax：baseline
  \item 差分注意力（DIFF）：差分降噪
  \item 高阶注意力（Nexus）：突破低秩瓶颈
\end{itemize}

现代 LLM 的注意力设计通常是从每个维度中
选择一个策略进行组合。
例如，DeepSeek-V3 选择了 MLA（维度一） + NSA（维度二），
Mistral 选择了 GQA（维度一） + SWA（维度二）。

\section{位置编码：让模型知道词的位置}
\label{sec:positional-encoding}

注意力机制有一个本质特性：
\textbf{它对输入的排列是不变的}（permutation equivariant）。
也就是说，如果你打乱输入词的顺序，
注意力的计算结果只是做了相应的排列：
因为点积只关心内容，不关心位置。

形式化地说，设 $\pi$ 是一个排列，
$P_\pi$ 是对应的置换矩阵，则：
\[
  \mathrm{Attention}(P_\pi Q, P_\pi K, P_\pi V)
  = P_\pi \, \mathrm{Attention}(Q, K, V)
\]
输入的顺序完全不影响每个 token 的注意力输出：
只是输出的排列跟着变了。

但语言显然是有顺序的。
「狗咬了人」和「人咬了狗」包含完全相同的词，
但意思截然相反。
模型必须知道每个词在句子中的位置。
位置编码正是为了打破注意力的排列不变性，
让模型能够区分不同的词序。

\subsection{从绝对到相对位置编码}

最初的 Transformer~\cite{vaswani2017attention} 使用正弦函数生成固定的位置编码，
直接加到词嵌入上：
\begin{align}
  \mathrm{PE}(pos, 2i) &= \sin\!\left(\frac{pos}{10000^{2i/d}}\right) \\
  \mathrm{PE}(pos, 2i+1) &= \cos\!\left(\frac{pos}{10000^{2i/d}}\right)
\end{align}

正弦编码的设计有其精妙之处：
不同维度对使用不同的频率（波长从 $2\pi$ 到 $10000 \cdot 2\pi$），
形成了一个类似\textbf{二进制编码}的多频率表示。
低维度对的高频正弦波区分相邻位置，
高维度对的低频正弦波编码粗粒度位置。
更关键的是，对于任意固定的偏移 $k$，
存在一个线性变换将 $\mathrm{PE}(pos)$ 映射到 $\mathrm{PE}(pos+k)$：
这意味着相对位置信息\emph{原则上}可以被线性层学习到。

但这种方法在实践中有两个问题：
一是模型只在有限长度上训练，无法自然泛化到更长的序列：
虽然正弦函数在数学上可以无限延伸，
但模型从未见过超出训练长度的位置编码模式，
因此无法正确解读它们；
二是绝对位置并不是最有用的信息：
对于语言理解来说，
「这两个词相距 3 个位置」（相对位置）
比「这个词在第 47 个位置」（绝对位置）更重要。

BERT 使用了可学习的绝对位置编码：
为每个位置学习一个独立的向量。
这比正弦编码更灵活，
但仍然是绝对位置，且不能泛化到训练时未见过的长度。

ALiBi（Attention with Linear Biases）~\cite{press2022alibi}
提出了另一种方案：
不修改嵌入，而是在注意力分数上直接加一个与距离成正比的偏置。
距离越远的 token 对，偏置越大（越负），注意力越低。
这自然编码了「近距离优先」的归纳偏置。

但最终胜出的是 RoPE。

\subsection{ALiBi：简洁的线性偏置方案}

在深入 RoPE 之前，值得更仔细地看看 ALiBi 的设计哲学，
因为它揭示了位置编码的一个关键洞察。

ALiBi 的做法极其简洁：在注意力分数 $QK^\top$ 上
加一个与 token 距离成线性正比的负偏置：
\begin{equation}
  \mathrm{Attention}(Q, K, V) =
  \mathrm{softmax}\!\left(\frac{QK^\top}{\sqrt{d_k}} + B\right) V
\end{equation}
其中偏置矩阵 $B_{ij} = -m \cdot |i - j|$，
$m$ 是每个注意力头特有的斜率参数
（头 $h$ 的斜率为 $2^{-8h/H}$，$H$ 是总头数）。

不同头使用不同的斜率意味着：
一些头（斜率小）几乎等权地关注整个序列，
另一些头（斜率大）几乎只关注邻近的 token。
这种\textbf{多尺度的位置感知}
无需任何可学习参数——斜率是固定的超参数。

ALiBi 的外推性能优于正弦编码和可学习绝对编码：
在 1K 上训练的模型可以在 2K 甚至 5K 上保持合理的困惑度，
而正弦编码在超出训练长度后迅速崩溃。
但 ALiBi 的线性衰减假设过于简单：
它隐含地认为注意力应该随距离单调递减，
而实际语言中存在大量的远距离依赖
（如首尾呼应的论述、远距离的指代关系）。
这限制了 ALiBi 在超长上下文场景中的表现。

\subsection{RoPE：旋转位置编码}

RoPE（Rotary Position Embedding）~\cite{su2021rope}
是当前所有主流 LLM 的标准位置编码方案。
它的核心思想优雅而巧妙：
\textbf{通过在复数空间中旋转向量来编码位置}。

\subsubsection{核心直觉}

为什么用「旋转」来编码位置？
考虑二维平面上的两个向量 $\mathbf{q}$ 和 $\mathbf{k}$。
如果我们分别将它们旋转角度 $m\theta$ 和 $n\theta$，
它们的点积变成：
\[
  \tilde{\mathbf{q}}^\top \tilde{\mathbf{k}}
  = \|\mathbf{q}\| \|\mathbf{k}\| \cos(\alpha + (m-n)\theta)
\]
其中 $\alpha$ 是原始夹角。
注意，结果只依赖于\textbf{相对位置} $m - n$，
而不是绝对位置 $m$ 或 $n$！
这正是我们想要的性质。

而且，旋转角度 $(m-n)\theta$ 意味着
相对距离越大，注意力分数的衰减越快：
这提供了一种自然的\textbf{远程衰减}（long-range decay）效应，
让模型更关注近距离的 token，
同时仍保留对远距离 token 的感知能力。

\subsubsection{数学公式}

对于位置 $m$ 的查询向量和位置 $n$ 的键向量，
RoPE 分别对它们施加旋转：
\begin{equation}
  \tilde{q}_m = R(\theta, m) \, q_m, \quad
  \tilde{k}_n = R(\theta, n) \, k_n
\end{equation}
其中 $R(\theta, m)$ 是一个旋转矩阵，
每对相邻维度 $(2i, 2i+1)$ 按角度 $m \theta_i$ 旋转：
\begin{equation}
  R(\theta_i, m) = \begin{pmatrix}
    \cos(m \theta_i) & -\sin(m \theta_i) \\
    \sin(m \theta_i) & \cos(m \theta_i)
  \end{pmatrix}, \quad
  \theta_i = 10000^{-2i/d}
\end{equation}

关键在于：旋转后的查询和键做点积时，
$\tilde{q}_m^\top \tilde{k}_n$ 只依赖于\textbf{相对位置 $m - n$}，
而不是绝对位置 $m$ 或 $n$。
这是因为旋转矩阵满足 $R(m)^\top R(n) = R(n - m)$。

$\theta_i$ 的设计也很有讲究：
低维度（$i$ 小）对应较大的 $\theta_i$，旋转速度快，
捕捉短距离的位置模式（如相邻词的语法关系）。
高维度（$i$ 大）对应较小的 $\theta_i$，旋转速度慢，
捕捉长距离的位置模式（如段落间的主题关联）。
这种多尺度编码使得模型能同时捕捉局部和全局的位置关系。

\subsubsection{旋转矩阵的几何本质}

RoPE 的优雅性最好通过几何来理解。
在二维平面上，旋转矩阵 $R(\theta)$ 将一个向量
绕原点旋转 $\theta$ 角：
这个操作保持向量的\textbf{模长不变}（$\|R(\theta)x\| = \|x\|$），
只改变方向。

高维的 RoPE 将 $d$ 维向量拆成 $d/2$ 个二维子平面，
每个子平面独立旋转。
整个旋转矩阵是一个\textbf{分块对角矩阵}：
\[
  R(\boldsymbol{\theta}, m) = \begin{pmatrix}
    R(\theta_1, m) & & \\
    & R(\theta_2, m) & \\
    & & \ddots \\
    & & & R(\theta_{d/2}, m)
  \end{pmatrix}
\]
每个 $2 \times 2$ 块以不同的角速度 $\theta_i$ 旋转。
这在物理上类似于\textbf{多频率的时钟}：
低维度的子平面像秒针（转得快，分辨短距离），
高维度的子平面像时针（转得慢，编码大范围位置）。

这种分块旋转结构有一个重要的工程优势：
它可以通过逐元素乘法和加法来高效实现，
而不需要显式构造和存储大旋转矩阵。
具体地，对于一个 $d$ 维向量 $x = (x_1, x_2, \ldots, x_d)$，
RoPE 的旋转操作等价于：
\[
  \tilde{x}_{2i} = x_{2i} \cos(m\theta_i) - x_{2i+1} \sin(m\theta_i), \quad
  \tilde{x}_{2i+1} = x_{2i} \sin(m\theta_i) + x_{2i+1} \cos(m\theta_i)
\]
这只需要 $2d$ 次乘法和 $d$ 次加法，
计算开销几乎可以忽略。

\subsubsection{为什么旋转优于加法？}

原始 Transformer 的正弦位置编码是\textbf{加法}的：
$\tilde{x} = x + \mathrm{PE}(m)$。
RoPE 是\textbf{乘法}的：$\tilde{x} = R(m) \cdot x$。

加法编码的问题在于位置信息和内容信息被混合在同一个向量空间中。
当两个向量做点积时：
\[
  (\mathbf{x}_m + \mathrm{PE}_m)^\top (\mathbf{x}_n + \mathrm{PE}_n)
  = \mathbf{x}_m^\top \mathbf{x}_n + \mathbf{x}_m^\top \mathrm{PE}_n
    + \mathrm{PE}_m^\top \mathbf{x}_n + \mathrm{PE}_m^\top \mathrm{PE}_n
\]
除了我们想要的纯位置项 $\mathrm{PE}_m^\top \mathrm{PE}_n$
和纯内容项 $\mathbf{x}_m^\top \mathbf{x}_n$ 之外，
还有两个\textbf{交叉项} $\mathbf{x}_m^\top \mathrm{PE}_n$
和 $\mathrm{PE}_m^\top \mathbf{x}_n$，
它们将位置和内容信息纠缠在一起，
使得模型难以干净地分离两者。

RoPE 通过乘法避免了这个问题：
\[
  (R_m q)^\top (R_n k) = q^\top R_m^\top R_n k = q^\top R(n-m) k
\]
结果只包含一个整体项，
其中位置信息 $R(n-m)$ 和内容信息 $q, k$ 以乘法形式耦合：
这种耦合方式更自然，
因为位置信息\emph{调制}了内容相似度，
而不是与之独立叠加。

\subsubsection{上下文长度外推}

RoPE 的一个实际挑战是\textbf{上下文长度外推}：
模型在 4K 长度上训练后，能否推广到 32K 或 128K？

直接外推通常会失败，因为模型没有见过高位置的旋转角度。
几种主流的解决方案：

\textbf{位置插值}（Position Interpolation）~\cite{chen2023extending}：
将长序列的位置缩放到训练时的范围。
例如，32K 的位置映射到 $[0, 4K)$，
相当于压缩了位置空间。
需要少量微调来适应压缩后的位置表示。

\textbf{NTK-aware 插值}~\cite{bloc97ntk}：
位置插值的问题在于它对所有频率维度做相同的缩放，
但不同频率维度的外推行为截然不同。
高频维度（旋转快的维度）对位置变化非常敏感：
即使微小的位置压缩也会改变它学到的局部模式；
低频维度（旋转慢的维度）本来就平滑，
对位置缩放更加鲁棒。

NTK-aware 插值的核心思想是修改 RoPE 的\textbf{基频}
$\theta_{\mathrm{base}}$（原始为 10000），
将其乘以一个扩展因子 $\alpha$：
\[
  \theta_i' = (\alpha \cdot 10000)^{-2i/d}
\]
当 $\alpha > 1$ 时，所有频率都被降低，
但\textbf{高频维度（$i$ 小）受影响最小}，
\textbf{低频维度（$i$ 大）受影响最大}。
这恰好是我们想要的：保留局部位置精度，
扩展全局位置范围。

NTK 名称来源于神经正切核（Neural Tangent Kernel）理论的类比：
NTK 理论中描述了高频成分难以学习的现象，
而 NTK-aware 插值正是利用了频率域中的
高/低频行为差异来做差异化处理。

\textbf{YaRN}（Yet another RoPE extensioN）进一步改进了 NTK-aware 插值，
引入了三项优化：
（1）\textbf{频率分段}，将维度分为高频、中频、低频三组，
对每组应用不同的缩放策略；
（2）\textbf{注意力温度缩放}——调整 softmax 的温度参数
来补偿频率修改对注意力分布的影响；
（3）\textbf{动态 NTK}——根据实际序列长度
自动调整缩放因子，短序列不缩放，
长序列逐渐增大缩放。

LLaMA 3 使用了修改后的 RoPE 频率
（$\theta_{\mathrm{base}} = 500{,}000$，远大于原始的 10{,}000），
配合特定维度的频率调整来支持 128K 上下文。
DeepSeek-V3 同样采用了 RoPE 频率调优方案，
并在 MLA 的位置编码分量中做了额外处理。

\begin{corebox}[RoPE 频率选择的工程经验]
RoPE 的 $\theta_{\mathrm{base}}$ 选择是一个重要的超参数。
较大的 $\theta_{\mathrm{base}}$（如 500K）使所有频率变慢，
有利于长上下文外推，但可能损害短距离位置精度。
较小的 $\theta_{\mathrm{base}}$（如 10K）
提供更好的短距离分辨率，但外推能力有限。
实践中的策略是：在预训练时使用中等 $\theta_{\mathrm{base}}$，
然后在长上下文微调阶段通过 NTK-aware 或 YaRN 方法扩展。
这种两阶段策略在 LLaMA 3 和 Qwen 2.5 中都被验证有效。
\end{corebox}

\section{前馈网络与激活函数}

Transformer 的每一层由两个子层组成：注意力层和前馈网络（FFN）层。
如果说注意力层负责「让词与词交流」，
那么 FFN 层则负责「在每个位置上独立地处理信息」。

这个「独立」很重要——FFN 层对序列中的每个位置
执行\textbf{完全相同}的变换，位置之间不共享信息。
从功能上看，FFN 层扮演着\textbf{知识存储}的角色。
Geva 等人~\cite{geva2021transformer} 的研究表明，
FFN 的第一层权重矩阵的每一行
可以解释为一个「模式检测器」：
激活特定的语义或语法模式，
第二层权重则将检测结果映射为对应的输出特征。
在这个视角下，FFN 本质上是一个\textbf{键值记忆}
（key-value memory），
其中键是第一层的行向量，值是第二层的列向量。

FFN 通常占据 Transformer 约\textbf{2/3 的参数量}
（在使用 SwiGLU 时约 67\%，标准 FFN 则恰好 2/3）。
这意味着大语言模型的大部分参数
实际上存储在「知识库」中，
而非「推理引擎」（注意力层）中。
这一观察对模型压缩和知识编辑有深远的启示。

\subsection{激活函数的演进}

激活函数的演进历史反映了深度学习社区对非线性建模的认知深化。

\textbf{ReLU}（2010s 主流）：$\mathrm{ReLU}(x) = \max(0, x)$。
计算简单、训练快，但存在「死亡 ReLU」问题：
一旦某个神经元的输入为负，梯度为零，该神经元永久失活。
原始 Transformer 的 FFN 就使用 ReLU：
$\mathrm{FFN}(x) = \max(0, xW_1 + b_1) W_2 + b_2$。

\textbf{GELU}（Gaussian Error Linear Unit, 2016）~\cite{hendrycks2016gelu}：
$\mathrm{GELU}(x) = x \cdot \Phi(x)$，
其中 $\Phi(x)$ 是标准正态分布的累积分布函数。
GELU 对负值不是硬截断，而是软性抑制，
提供了更平滑的梯度。BERT 和 GPT-2 使用了 GELU。

\textbf{SiLU/Swish}（2017）~\cite{ramachandran2017swish}：
$\mathrm{Swish}(x) = x \cdot \sigma(\beta x)$，
其中 $\sigma$ 是 sigmoid 函数，$\beta$ 通常为 1。
Swish 在数学上与 GELU 非常接近，
但形式更简洁。Google Brain 通过自动搜索发现了这个激活函数。

\textbf{SwiGLU}（2020）~\cite{shazeer2020glu}：
将 Swish 与门控线性单元（GLU）结合：
\begin{equation}
  \mathrm{SwiGLU}(x) = \left(\mathrm{Swish}(x W_1) \odot (x W_3)\right) W_2
\end{equation}
其中 $\mathrm{Swish}(x) = x \cdot \sigma(x)$，
$\odot$ 是逐元素乘法。SwiGLU 引入了一个门控机制：
$xW_3$ 控制哪些信息被通过，
$\mathrm{Swish}(xW_1)$ 提供非线性变换。

\subsection{GLU 家族：门控线性单元的统一视角}

在理解 SwiGLU 之前，有必要看看完整的 GLU 家族。
GLU（Gated Linear Unit）由 Dauphin 等人~\cite{dauphin2017glu}
在 2017 年提出，原始形式为：
\begin{equation}
  \mathrm{GLU}(x) = (x W_1 + b_1) \otimes \sigma(x W_3 + b_3)
\end{equation}
其中 $\sigma$ 是 sigmoid 函数，$\otimes$ 是逐元素乘法。

GLU 的设计灵感来自 LSTM：
$\sigma(xW_3 + b_3)$ 输出 $[0,1]$ 之间的值，
作为「门」控制 $xW_1 + b_1$ 中每个维度的信息是否通过。
当门接近 0 时，信息被完全阻断；
当门接近 1 时，信息完全通过。

替换不同的激活函数，就得到了 GLU 家族的不同成员：

\begin{center}\small
\renewcommand{\arraystretch}{1.3}
\begin{tabular}{lll}
\toprule
\rowcolor{figLightGray}
\textbf{变体} & \textbf{公式} & \textbf{代表模型} \\
\midrule
GLU & $(xW_1) \otimes \sigma(xW_3)$ & 原始论文 \\
ReGLU & $\mathrm{ReLU}(xW_1) \otimes (xW_3)$ & --- \\
GEGLU & $\mathrm{GELU}(xW_1) \otimes (xW_3)$ & Gemma \\
SwiGLU & $\mathrm{Swish}(xW_1) \otimes (xW_3)$ & LLaMA, PaLM \\
\bottomrule
\end{tabular}
\end{center}

SwiGLU 成为事实标准并非偶然：
Shazeer~\cite{shazeer2020glu} 在系统性的对比实验中发现，
SwiGLU 和 GEGLU 在所有基准上一致优于其他变体，
而两者之间的差异极小。
LLaMA 选择 SwiGLU，Gemma 选择 GEGLU，
更多是工程习惯而非性能差异。

\subsection{为什么门控机制有效？}

门控机制的核心思想是\textbf{信息过滤}。
在标准 FFN 中，所有信息经过相同的非线性变换。
在门控 FFN 中，模型学习了一个「开关」：
决定哪些维度的信息应该通过，哪些应该被抑制。
这类似于 LSTM 中的遗忘门和输入门。

直觉上，语言模型在每个位置上需要做两件事：
（1）提取有用的特征，（2）过滤无关的噪声。
门控机制让模型显式地学习这两个步骤，
而不是将它们混在一起。

从梯度流的角度看，门控机制还有一个隐蔽的优势。
在标准 ReLU FFN 中，$\mathrm{ReLU}(x) = \max(0, x)$
的梯度在 $x < 0$ 时为零：
这意味着负半轴的梯度完全消失，
该路径上的参数得不到更新，
即所谓的「死亡 ReLU」问题。
在 SwiGLU 中，即使 $\mathrm{Swish}(xW_1)$ 的某些维度很小，
梯度仍然可以通过 $xW_3$ 分支反向传播：
门控路径提供了\textbf{备用梯度通道}，
缓解了梯度消失问题。

Shazeer~\cite{shazeer2020glu} 的实验表明，
在控制参数量相同的条件下，
SwiGLU 在多个基准上一致优于 ReLU 和 GELU。
这个优势虽然不大（通常 1--3 个百分点），
但在大规模训练中累积起来就是显著的差异。
当训练预算以百万美元计时，
每个百分点的效率提升都有巨大的经济意义。

\begin{whybox}[为什么 FFN 的隐藏维度是 $8d/3$？]
原始 Transformer 中 FFN 的隐藏维度是 $4d$（$d$ 是模型维度）。
FFN 有两个矩阵 $W_1 \in \mathbb{R}^{d \times 4d}$ 和 $W_2 \in \mathbb{R}^{4d \times d}$，
总参数量为 $2 \times 4d^2 = 8d^2$。
SwiGLU 引入了第三个权重矩阵 $W_3$，增加了约 50\% 的参数。
为了保持总参数量不变（$3 \times d \times d_{\mathrm{ffn}} = 8d^2$），
隐藏维度从 $4d$ 缩小到 $d_{\mathrm{ffn}} = 8d/3 \approx 2.67d$。
这样 SwiGLU 在相同参数预算下，
性能优于更大的 ReLU FFN。
实践中，这个值通常被取整到 256 的倍数
以便于 GPU tensor core 对齐，
例如 LLaMA 3-8B 使用 $d = 4096, d_{\mathrm{ffn}} = 14336$
（恰好是 $3.5d$，略大于 $8d/3$）。
\end{whybox}

\section{层归一化：训练的稳定器}

深层神经网络的训练面临一个根本挑战：
随着层数增加，各层的输入分布可能剧烈变化，
导致训练不稳定甚至发散。
这个现象被称为\textbf{内部协变量偏移}
（Internal Covariate Shift）~\cite{ioffe2015batchnorm}：
虽然这个解释在学术上存在争议，
但归一化技术在实践中的有效性毋庸置疑。

层归一化（Layer Normalization）~\cite{ba2016layernorm}
通过对每一层的输入做标准化来缓解这个问题。
与批归一化（Batch Normalization）不同，
层归一化在\textbf{特征维度}而非批次维度上计算统计量：
这使得它不依赖于 batch size，
适合自回归语言模型的变长序列处理。

现代 LLM 在层归一化上做了两个关键改进。

\subsection{Pre-Norm vs Post-Norm}

原始 Transformer 在注意力/FFN \textbf{之后}做归一化（Post-Norm）：
\[
  x' = \mathrm{Norm}(x + \mathrm{SubLayer}(x))
\]
现代模型在\textbf{之前}做归一化（Pre-Norm）：
\[
  x' = x + \mathrm{SubLayer}(\mathrm{Norm}(x))
\]

这两种位置看似只是微小的差异，但影响深远。

\textbf{Pre-Norm 的优势}在于残差连接的梯度流。
在 Pre-Norm 中，残差连接形成了一条从最后一层到第一层的
「梯度高速公路」——$x' = x + f(\mathrm{Norm}(x))$，
其中 $\partial x'/\partial x = I + \partial f / \partial x$。
即使 $\partial f / \partial x$ 很小或为零，
梯度也能通过恒等映射 $I$ 畅通无阻地反向传播。
这使得训练 100+ 层的深层模型成为可能。

\textbf{Post-Norm 的优势}在于归一化后的表示更稳定。
由于归一化在残差连接之后，
每一层的输出都被归一化到统一尺度，
理论上表示能力更强。
但代价是深层模型容易出现梯度消失或训练不稳定。

实践中，几乎所有现代 LLM 都选择 Pre-Norm：
训练稳定性比微小的表示能力差异更重要。
DeepSeek-V3 在 Pre-Norm 基础上进一步引入了
额外的辅助损失来防止训练崩溃。

\subsection{梯度流的定量分析}

Pre-Norm 和 Post-Norm 的差异可以通过梯度的传播路径精确量化。

考虑一个 $L$ 层的 Transformer。
在 Pre-Norm 中，第 $l$ 层的前向传播为：
$x_{l+1} = x_l + f_l(\mathrm{Norm}(x_l))$。
从第 $L$ 层反传到第 $l$ 层的梯度为：
\[
  \frac{\partial x_L}{\partial x_l} = I + \sum_{k=l}^{L-1}
  \frac{\partial f_k}{\partial x_l} \cdot \prod_{j=k+1}^{L-1}
  \left(I + \frac{\partial f_j}{\partial x_j}\right)
\]
关键项是\textbf{恒等矩阵 $I$}：
无论 $f_k$ 的梯度如何，$I$ 都保证了至少有一条
从第 $L$ 层到第 $l$ 层的\textbf{无衰减梯度路径}。
这就是所谓的「梯度高速公路」。

在 Post-Norm 中，第 $l$ 层的前向传播为：
$x_{l+1} = \mathrm{Norm}(x_l + f_l(x_l))$。
由于归一化操作包含除法，
梯度路径中不再有简单的 $I$ 项：
每一层的归一化都会对梯度做非线性变换，
梯度在深层网络中可能指数级衰减或爆炸。

Xiong 等人~\cite{xiong2020prenorm} 严格证明了：
在 Post-Norm Transformer 中，
如果不精心调整学习率预热（warmup），
梯度范数在训练初期可能在层间呈指数增长，
导致优化器收到错误的更新方向。
Pre-Norm 则完全避免了这个问题：
即使没有 warmup，训练也能稳定进行。

\subsection{DeepNorm：突破深度限制}

Pre-Norm 让训练稳定了，但随着层数进一步增加到数百甚至上千层，
即使 Pre-Norm 也面临挑战：
残差连接的累积可能导致隐状态的范数逐层增长，
最终造成数值问题。

\textbf{DeepNorm}~\cite{wang2022deepnet} 是 Microsoft Research 提出的
针对超深 Transformer 的归一化方案。
核心修改只有一个公式的微调：
\[
  x_{l+1} = \mathrm{Norm}(\alpha \cdot x_l + f_l(x_l))
\]
其中 $\alpha > 1$ 是一个与层数相关的缩放因子
（$\alpha = (2L)^{1/4}$），
配合对残差分支权重的 Xavier 初始化缩放
（$\beta = (8L)^{-1/4}$）。

$\alpha > 1$ 的作用是\textbf{放大残差路径相对于子层路径}的权重，
使得信号的主要成分通过残差通道传递，
子层的贡献被相对抑制：
这在深层网络中防止了信号被逐层累积的子层输出所主导。

DeepNorm 使得训练 1000 层的 Transformer 成为可能
（而标准 Pre-Norm 在约 100 层后开始出现稳定性问题），
并且在相同参数量下，更深的模型可以获得比更宽模型更好的性能。

\subsection{RMSNorm：去掉均值，只留尺度}

标准 LayerNorm 的计算过程包括两步：
均值平移（center）和尺度归一化（scale）：
\begin{equation}
  \mathrm{LayerNorm}(x) = \frac{x - \mu}{\sqrt{\sigma^2 + \epsilon}} \cdot \gamma + \beta
\end{equation}
其中 $\mu = \frac{1}{d}\sum_{i} x_i$ 是均值，
$\sigma^2 = \frac{1}{d}\sum_{i}(x_i - \mu)^2$ 是方差，
$\gamma$ 和 $\beta$ 是可学习的缩放和偏移参数。

\textbf{RMSNorm}~\cite{zhang2019rms} 去掉了均值平移操作，
只做尺度归一化：
\begin{equation}
  \mathrm{RMSNorm}(x) = \frac{x}{\sqrt{\frac{1}{d}\sum_{i=1}^{d} x_i^2 + \epsilon}} \cdot \gamma
\end{equation}

从数学上分析，LayerNorm 的均值平移操作
相当于将向量投影到垂直于全 1 向量 $\mathbf{1}$ 的超平面上。
Zhang 和 Sennrich~\cite{zhang2019rms} 的研究表明，
\textbf{归一化的主要贡献来自尺度不变性，而非均值平移}。
去掉均值平移不影响模型质量，
但节省了计算均值和减去均值的操作：
在大规模模型中，这带来了约 10--15\% 的归一化层加速。
RMSNorm 还去掉了偏移参数 $\beta$，进一步简化了模型。

从直觉上理解 RMSNorm 的有效性：
LayerNorm 的均值平移让向量的均值为零，
RMSNorm 不做这一步——意味着允许向量有非零均值。
在高维空间中，向量的均值成分
（所有维度的平均值）只是一个维度的信息，
相对于 $d$ 维的总信息量微不足道。
真正重要的是\textbf{尺度归一化}：
确保不同层、不同位置的向量在数值上处于可比的范围。

RMSNorm 已成为 2024--2026 年所有主流 LLM 的标准选择，
包括 LLaMA 3、Mistral、DeepSeek-V3、Qwen 2.5、Gemma 2 等。
没有任何 frontier 模型使用标准 LayerNorm：
RMSNorm 是一个罕见的「只有好处没有代价」的改进。

\section{Flash Attention：让注意力更快}
\label{sec:flash-attention}

注意力机制的计算复杂度是 $O(n^2)$：
对于长序列，这既是计算瓶颈，也是内存瓶颈。
标准实现需要在 GPU 的高带宽内存（HBM）中存储完整的 $n \times n$ 注意力矩阵。

Flash Attention 的创造者 Tri Dao
在 Stanford PhD 期间的核心洞察是：
深度学习社区长期以来用 FLOPS（浮点运算数）衡量算法效率，
但在现代 GPU 上，\textbf{内存传输}才是真正的瓶颈。
一个算法即使做了更少的浮点运算，
如果频繁访问慢速内存，实际运行时间可能更长。
这就是 \textbf{IO-awareness}（IO 感知）的核心理念：
算法设计必须考虑硬件的内存层次结构。

\subsection{GPU 内存层次结构}

要理解 Flash Attention 为什么有效，
必须先理解 GPU 的内存层次：

\begin{center}\small
\renewcommand{\arraystretch}{1.3}
\begin{tabular}{lcc}
\toprule
\rowcolor{figLightGray}
\textbf{层级} & \textbf{容量} & \textbf{带宽} \\
\midrule
寄存器（Register） & $\sim$KB / SM & $\sim$TB/s \\
SRAM（片上缓存） & 20\,MB（A100） & $\sim$19\,TB/s \\
HBM（高带宽内存） & 80\,GB（A100） & $\sim$2\,TB/s \\
\bottomrule
\end{tabular}
\end{center}

SRAM 的带宽是 HBM 的约 \textbf{10 倍}，
但容量只有 HBM 的约 \textbf{1/4000}。
标准注意力实现将 $Q, K, V$ 和 $n \times n$ 注意力矩阵都存在 HBM 中，
计算过程不断在 HBM 和 SRAM 之间搬运数据，
而 GPU 的算术单元大部分时间在等待数据传输。
这就是为什么标准注意力是 \textbf{memory-bound}（内存带宽受限）
而非 compute-bound（计算受限）的。

\subsection{IO 复杂度分析}

为了量化内存瓶颈，Dao 等人引入了
\textbf{IO 复杂度}——衡量算法在 HBM 上的读写次数。

标准注意力的实现需要以下 HBM 操作：
\begin{enumerate}[noitemsep]
  \item 从 HBM 读取 $Q, K$（各 $O(nd)$ 元素）
  \item 将 $QK^\top$ 的 $n \times n$ 结果写回 HBM（$O(n^2)$ 元素）
  \item 从 HBM 读取 $n \times n$ 矩阵做 softmax，写回结果
  \item 从 HBM 读取 softmax 结果和 $V$，计算最终输出
\end{enumerate}
总 IO 复杂度为 $O(nd + n^2)$：
当 $n > d$ 时，$n^2$ 项主导，
这意味着大部分时间花在了读写那个巨大的注意力矩阵上。

Flash Attention 的目标是将 IO 复杂度降至 $O(n^2 d^2 / M)$，
其中 $M$ 是 SRAM 的大小。
当 $M = O(nd)$（SRAM 能容纳 $Q$ 的几行和 $K, V$ 的几行）时，
IO 复杂度降为 $O(nd)$——\textbf{线性于序列长度}，
而非标准实现的 $O(n^2)$。

更重要的是，Dao 证明了这个 IO 复杂度是\textbf{最优的}：
在给定的 SRAM 大小 $M$ 下，不存在渐近更好的算法。
这意味着 Flash Attention 不仅是一个好的实现，
而且是\textbf{理论上最优的}注意力 IO 算法。

\subsection{核心算法：分块计算}

Flash Attention~\cite{dao2022flashattention,dao2023flashattention2}
的核心洞察是：
\textbf{注意力计算是内存带宽瓶颈的（memory-bound），而非计算瓶颈的}。
GPU 的算术运算速度远快于内存读写速度。
标准实现在 HBM 和 SRAM 之间来回搬运数据的时间，
远超实际的矩阵运算时间。

Flash Attention 的解决方案是\textbf{分块计算（tiling）}：
\begin{enumerate}[noitemsep]
  \item 将 $Q$ 分成块 $Q_1, Q_2, \ldots$，将 $K, V$ 也分成对应的块
  \item 对于每个 $Q$ 块，依次加载 $K, V$ 的块到 SRAM
  \item 在 SRAM 中完成 $Q_i K_j^\top$ 的计算、softmax、以及与 $V_j$ 的乘法
  \item 使用\textbf{在线 softmax}（online softmax）算法
        增量地更新输出：
        不需要先计算整行的注意力分数再做 softmax，
        而是在处理每个 $K, V$ 块时动态更新 softmax 的分母
  \item 最终结果直接写回 HBM，中间的 $n \times n$ 注意力矩阵\textbf{从不存储}
\end{enumerate}

结果是：Flash Attention 在数学上\textbf{完全等价}于标准注意力
（不是近似！），但内存占用从 $O(n^2)$ 降至 $O(n)$，
速度提升 2--4$\times$。

\subsection{分块算法的伪代码}

为了让分块计算的过程完全透明，
以下是 Flash Attention 前向传播的简化伪代码：

\begin{lstlisting}[caption={Flash Attention forward pass (simplified)}]
def flash_attention_forward(Q, K, V, block_size_q, block_size_kv):
    """
    Q: (n, d_k), K: (n, d_k), V: (n, d_v)
    Returns: O (n, d_v) — same result as standard attention
    """
    n, d = Q.shape
    O = zeros(n, d)       # output accumulator
    l = zeros(n)          # softmax denominator accumulator
    m = full(n, -inf)     # running max for numerical stability

    # Outer loop: iterate over Q blocks
    for i in range(0, n, block_size_q):
        Qi = Q[i:i+block_size_q]        # load Q block to SRAM
        Oi = zeros(block_size_q, d)     # local output
        li = zeros(block_size_q)        # local denominator
        mi = full(block_size_q, -inf)   # local max

        # Inner loop: iterate over K, V blocks
        for j in range(0, n, block_size_kv):
            Kj = K[j:j+block_size_kv]   # load K block to SRAM
            Vj = V[j:j+block_size_kv]   # load V block to SRAM

            # Compute block attention scores
            Sij = Qi @ Kj.T / sqrt(d)   # (block_q, block_kv)

            # Online softmax update
            mij = Sij.max(dim=-1)
            mi_new = max(mi, mij)
            Pij = exp(Sij - mi_new[:, None])
            li_new = li * exp(mi - mi_new) + Pij.sum(dim=-1)
            Oi = Oi * (li * exp(mi - mi_new))[:, None] \
                 + Pij @ Vj
            mi, li = mi_new, li_new

        # Normalize and write back to HBM
        O[i:i+block_size_q] = Oi / li[:, None]

    return O
\end{lstlisting}

这段伪代码的关键在于\textbf{外循环遍历 Q 块，内循环遍历 KV 块}。
对于每个 Q 块，所有 KV 块被逐一加载到 SRAM 中处理，
在线 softmax 增量更新输出。
整个过程中，\textbf{$n \times n$ 的注意力矩阵从未被完整存储}：
每次只在 SRAM 中存储一个
$\mathrm{block\_q} \times \mathrm{block\_kv}$ 大小的块，
处理完后立即丢弃。

\begin{corebox}[在线 softmax 的关键]
标准 softmax 需要两遍扫描：第一遍求分母 $\sum_j e^{s_j}$，
第二遍归一化。
在线 softmax 在一遍扫描中完成：
维护一个运行中的最大值 $m$ 和分母 $l$，
每处理一个新块时更新：
\[
  m_{\mathrm{new}} = \max(m_{\mathrm{old}}, m_{\mathrm{block}}), \quad
  l_{\mathrm{new}} = l_{\mathrm{old}} \cdot e^{m_{\mathrm{old}} - m_{\mathrm{new}}}
  + l_{\mathrm{block}} \cdot e^{m_{\mathrm{block}} - m_{\mathrm{new}}}
\]
这使得分块计算在数学上精确等价于一次性计算。
\end{corebox}

\subsection{Flash Attention 的演进}

\textbf{Flash Attention 1}~\cite{dao2022flashattention}（2022）
引入了基本的分块算法和 IO-awareness 概念。
在 A100 GPU 上实现了 2--4$\times$ 的加速和显著的内存节省。

\textbf{Flash Attention 2}~\cite{dao2023flashattention2}（2023）
在算法层面做了优化：
\begin{itemize}[noitemsep]
  \item 减少了非矩阵乘法操作的数量（这些操作在 GPU 上效率较低）
  \item 改进了 GPU 线程块和 warp 之间的并行化策略
  \item 将序列长度维度的并行化改为在注意力头和 batch 维度上并行
  \item 在 A100 上达到了理论峰值 FLOPS 的 50--73\%，
        而 Flash Attention 1 仅达到 25--40\%
\end{itemize}

\textbf{Flash Attention 3}（2024）~\cite{shah2024flashattention3}
针对 H100 GPU（Hopper 架构）的新特性做了深度优化。
Flash Attention 2 的设计面向 Ampere 架构（A100），
在 H100 上仅能达到约 35\% 的利用率：
远未充分利用 Hopper 的新硬件能力。
Flash Attention 3 通过三项关键技术释放了 Hopper 的潜力：

\textbf{（1）Warp 专业化与异步流水线}。
Hopper 架构引入了两个关键新特性：
\textbf{WGMMA}（Warp Group Matrix Multiply Accumulate）指令
可以直接从共享内存执行矩阵乘法，
\textbf{TMA}（Tensor Memory Accelerator）
可以异步地在 HBM 和共享内存之间搬运数据。
Flash Attention 3 将 warp 分为\textbf{生产者}和\textbf{消费者}两组：
生产者 warp 使用 TMA 异步加载下一个数据块到共享内存，
消费者 warp 使用 WGMMA 对当前数据块执行矩阵乘法。
通过环形共享内存缓冲区，
数据加载和计算可以完全\textbf{重叠}：
GPU 不再有任何空闲周期等待数据。

\textbf{（2）两阶段 GEMM-softmax 流水线（Ping-Pong 调度）}。
注意力计算中有两次矩阵乘法：$QK^\top$ 和 $\mathrm{softmax} \times V$。
传统实现必须等 $QK^\top$ 完成后才能开始 softmax，
再等 softmax 完成后才能执行 $\times V$。
Flash Attention 3 将相邻迭代的操作交错执行：
在第 $j+1$ 次迭代的 $QK^\top$ 计算同时，
执行第 $j$ 次迭代的 softmax。
这种 ping-pong 调度进一步隐藏了非矩阵乘法操作的延迟。

\textbf{（3）FP8 块量化与不相干处理}。
H100 的 FP8 tensor core 提供高达 \textbf{1.98 PFLOPS} 的理论峰值。
Flash Attention 3 利用\textbf{块量化}（block quantization）
将 FP16 的 $Q, K$ 在块级别转换为 FP8，
配合\textbf{不相干处理}（incoherent processing）：
通过随机正交变换使量化误差分布更均匀。
结果是 FP8 Flash Attention 3 的数值误差
比 baseline FP8 注意力低 \textbf{2.6$\times$}，
吞吐量接近 \textbf{1.2 PFLOPS}。

\begin{center}\small
\renewcommand{\arraystretch}{1.3}
\begin{tabular}{lcccc}
\toprule
\rowcolor{figLightGray}
\textbf{版本} & \textbf{GPU} & \textbf{精度} & \textbf{吞吐量} & \textbf{峰值利用率} \\
\midrule
标准注意力 & A100 & FP16 & --- & $<$20\% \\
Flash Attention 1 & A100 & FP16 & --- & 25--40\% \\
Flash Attention 2 & A100 & FP16 & --- & 50--73\% \\
Flash Attention 3 & H100 & FP16 & 740 TFLOPS & $\sim$75\% \\
Flash Attention 3 & H100 & FP8 & $\sim$1.2 PFLOPS & $\sim$60\% \\
\bottomrule
\end{tabular}
\end{center}

Flash Attention 3 使得在 H100 上处理超长上下文变得实际可行：
740 TFLOPS 的 FP16 吞吐量意味着
处理 128K 上下文的注意力计算时间
比 Flash Attention 2 在 A100 上快约 3--4$\times$
（考虑 H100 本身的硬件提升）。
FP8 模式更是将吞吐量推到 PFLOPS 级别，
为 1M+ token 的上下文处理开辟了道路。

截至 2026 年初，尚未有官方的 “Flash Attention 4” 发布，
但 Flash Attention 3 已在 GTC 2025 上
被 Tri Dao 演示了面向 Blackwell（B200）架构的优化方向：
包括利用 B200 的第五代 Tensor Core（FP4 支持）
和更大的 SRAM（每 SM 提升至 32\,MB 级别）。

Flash Attention 的成功故事揭示了一个深层教训：
\textbf{算法创新不一定意味着降低计算复杂度}。
Flash Attention 的计算量与标准注意力\emph{完全相同}
（都是 $O(n^2 d)$ FLOPs），
它的加速完全来自\textbf{减少内存传输}。
在 GPU 时代，内存带宽而非计算能力才是真正的瓶颈：
理解这一点对于设计高效的深度学习算法至关重要。

\textbf{Flash Attention 3 的生产部署}（2026 年 3 月）
标志着 FA3 从论文进入真正的工业级应用。
在 Blackwell GPU 上，FA3 通过 \textbf{cuTile} 实现了
面向生产环境的高效注意力内核，
引入了多项工程优化：
FMA（fused multiply-add）模式减少了指令数量，
fast math 路径在精度可接受的场景下进一步提升吞吐量，
loop splitting 和 adaptive tiling 根据实际序列长度
动态调整分块策略以最大化硬件利用率。

在 PyTorch 生态中，FA3 已作为 \texttt{third\_party} 子模块
正式集成到主仓库中。
值得注意的是，FA3 的硬件支持已不再局限于 H100：
通过适配层实现了对 A100（Ampere 架构）的支持，
使得更广泛的 GPU 集群可以受益于 FA3 的优化。

此外，针对 DeepSeek 的 MLA 架构，
社区开发了 \textbf{FlashMLA} 实现：
在 Flash Attention 的 IO-aware 框架上
专门优化了 MLA 的低秩 KV 投影计算路径，
使得 MLA 在推理中的实际速度进一步提升。

Flash Attention 已经成为所有现代 LLM 训练和推理的基础设施，
集成在 PyTorch 2.0+、Hugging Face Transformers、
以及几乎所有主流训练框架中。
值得注意的是，Flash Attention 的 IO-awareness 理念
也被应用到了 SSM 领域：
Mamba-2 的 SSD 算法就采用了类似的分块策略
来最大化 GPU 硬件利用率。

\section{MLA：DeepSeek 的注意力压缩}

DeepSeek-V2 引入了\textbf{多头潜在注意力}
（Multi-head Latent Attention, MLA），
进一步解决 KV cache 的内存瓶颈。

图~\ref{fig:kv-sharing}直观展示了
从 MHA 到 GQA 到 MQA 再到 MLA 的 KV 共享策略演进。

\begin{figure}[htbp]
\centering
\begin{tikzpicture}[
  qhead/.style={draw=figBlue!50, fill=figBlue!10, rounded corners=1.5pt,
    minimum width=8mm, minimum height=11mm, font=\scriptsize, inner sep=1.5pt,
    line width=0.6pt},
  kvhead/.style={draw=figOrange!50, fill=figOrange!10, rounded corners=1.5pt,
    minimum height=11mm, font=\scriptsize, inner sep=2pt, line width=0.6pt},
  latent/.style={draw=figGreen!50, fill=figGreen!10, rounded corners=1.5pt,
    minimum height=11mm, font=\scriptsize, inner sep=3pt, line width=0.6pt},
  title/.style={font=\small\bfseries, text=black!80},
  sub/.style={font=\tiny, text=figGray},
  conn/.style={-{Stealth[length=2.5pt,width=2.5pt]}, line width=0.6pt, figGray!60},
  progression/.style={-{Stealth[length=2.5pt,width=2.5pt]}, line width=0.6pt,
    figGray!40, dashed},
]

\def\colsep{40mm}

% === MHA ===
\begin{scope}
  \node[title] at (0, 3.2) {MHA};
  \foreach \i in {1,...,4} {
    \node[qhead] (mha-q\i) at ({(\i-2.5)*1.05}, 1.8) {\textbf{Q}$_\i$};
    \node[kvhead, minimum width=8mm] (mha-kv\i) at ({(\i-2.5)*1.05}, 0) {\textbf{KV}$_\i$};
    \draw[conn] (mha-q\i) -- (mha-kv\i);
  }
  \node[sub] at (0, -1.0) {Cache: $O(h \cdot d)$};
\end{scope}

% === GQA ===
\begin{scope}[xshift=\colsep]
  \node[title] at (0, 3.2) {GQA};
  \foreach \i in {1,...,4} {
    \node[qhead] (gqa-q\i) at ({(\i-2.5)*1.05}, 1.8) {\textbf{Q}$_\i$};
  }
  \node[kvhead, minimum width=16mm] (gqa-kv1) at (-1.05, 0) {\textbf{KV}$_1$};
  \node[kvhead, minimum width=16mm] (gqa-kv2) at (1.05, 0) {\textbf{KV}$_2$};
  \draw[conn] (gqa-q1) -- (gqa-kv1);
  \draw[conn] (gqa-q2) -- (gqa-kv1);
  \draw[conn] (gqa-q3) -- (gqa-kv2);
  \draw[conn] (gqa-q4) -- (gqa-kv2);
  \node[sub] at (0, -1.0) {Cache: $O(g \cdot d)$, $g < h$};
\end{scope}

% === MQA ===
\begin{scope}[xshift=2*\colsep]
  \node[title] at (0, 3.2) {MQA};
  \foreach \i in {1,...,4} {
    \node[qhead] (mqa-q\i) at ({(\i-2.5)*1.05}, 1.8) {\textbf{Q}$_\i$};
  }
  \node[kvhead, minimum width=32mm] (mqa-kv) at (0, 0) {\textbf{KV} (shared)};
  \foreach \i in {1,...,4} {
    \draw[conn] (mqa-q\i) -- (mqa-kv);
  }
  \node[sub] at (0, -1.0) {Cache: $O(d)$};
\end{scope}

% === MLA ===
\begin{scope}[xshift=3*\colsep]
  \node[title] at (0, 3.2) {MLA};
  \foreach \i in {1,...,4} {
    \node[qhead] (mla-q\i) at ({(\i-2.5)*1.05}, 1.8) {\textbf{Q}$_\i$};
  }
  \node[latent, minimum width=32mm] (mla-c) at (0, 0)
    {Latent $\mathbf{c}$\enspace($d_c \!\ll\! d$)};
  \foreach \i in {1,...,4} {
    \draw[conn] (mla-q\i) -- (mla-c);
  }
  \node[sub] at (0, -1.0) {Cache: $O(d_c)$};
\end{scope}

% === Progression arrows ===
\draw[progression] (0.9*\colsep, 1.0) -- (1.1*\colsep, 1.0);
\draw[progression] (1.9*\colsep, 1.0) -- (2.1*\colsep, 1.0);
\draw[progression] (2.9*\colsep, 1.0) -- (3.1*\colsep, 1.0);

\end{tikzpicture}
\caption{KV 共享策略的演进。MHA 中每个查询头拥有独立 KV 头（$h$ 份缓存）；
GQA 将多个查询头映射到同一 KV 组；MQA 进一步共享为单一 KV 头；
MLA（DeepSeek-V2）将全部 KV 压缩为低维潜在向量 $\mathbf{c}$，
缓存开销从 $O(hd)$ 降至 $O(d_c)$。}
\label{fig:kv-sharing}
\end{figure}

\subsection{动机与设计思想}

GQA 通过共享 KV 减少了 KV 头的数量，
但每个 KV 头的维度不变。
MLA 的思路不同：\textbf{将 KV 投影到低维潜在空间}。

核心观察是：在高维空间中，
$K$ 和 $V$ 矩阵存在大量冗余：
它们的有效秩（effective rank）远低于矩阵维度。
MLA 利用这个低秩结构，
将 KV 压缩到一个低维的\textbf{潜在向量} $c$。

\subsection{数学形式}

具体而言，MLA 不直接缓存 $K$ 和 $V$，
而是缓存一个压缩的\textbf{潜在向量} $c$（维度远小于 $K, V$），
在计算注意力时再将 $c$ 投影回 $K, V$：
\begin{equation}
  c = W_{\mathrm{down}} \, x, \quad
  K = W_K^{\mathrm{up}} \, c, \quad
  V = W_V^{\mathrm{up}} \, c
\end{equation}
其中 $W_{\mathrm{down}} \in \mathbb{R}^{d_c \times d}$（$d_c \ll d$）
是下投影矩阵。

关键的工程细节是\textbf{矩阵吸收}（absorption trick）：
在推理时，注意力分数的计算为：
\[
  q^\top K = q^\top W_K^{\mathrm{up}} c
  = (W_K^{\mathrm{up}\top} q)^\top c
  = \tilde{q}^\top c
\]
其中 $\tilde{q} = W_K^{\mathrm{up}\top} q$ 是一个\textbf{预变换的查询}，
维度为 $d_c$（远小于原始的 $d$）。
这意味着注意力可以\textbf{直接在低维空间 $d_c$ 中计算}：
不需要将 $c$ 展开为完整的 $K$ 矩阵，
而是将上投影矩阵 $W_K^{\mathrm{up}}$ 吸收到查询的投影中。

同样的技巧应用于值的投影：
$\mathrm{Output} = A \cdot V = A \cdot W_V^{\mathrm{up}} c$，
输出投影矩阵可以吸收 $W_V^{\mathrm{up}}$，
使得整个注意力层在推理时只操作 $d_c$ 维的向量。

这个吸收技巧是 MLA 的工程精髓：
它使得 MLA 不仅在\emph{存储}上高效
（只缓存低维的 $c$），
在\emph{计算}上也同样高效
（注意力在低维空间中完成）。

\subsection{压缩效果}

通过将 KV 压缩到 $d_c$ 维（DeepSeek-V2 中 $d_c = 512$，
而原始 KV 维度为 $n_h \times d_h = 128 \times 128 = 16{,}384$），
MLA 将 KV cache 减少到 GQA 的约 $1/6$，
同时保持与 MHA 相当的质量。

\begin{center}\small
\renewcommand{\arraystretch}{1.3}
\begin{tabular}{lcc}
\toprule
\rowcolor{figLightGray}
\textbf{方法} & \textbf{每 token KV cache 大小} & \textbf{相对大小} \\
\midrule
MHA（64 头） & $2 \times 64 \times 128 = 16{,}384$ & $1\times$ \\
GQA（8 KV 头） & $2 \times 8 \times 128 = 2{,}048$ & $1/8$ \\
MQA（1 KV 头） & $2 \times 1 \times 128 = 256$ & $1/64$ \\
MLA（$d_c = 512$） & $512$ & $\sim 1/32$ \\
\bottomrule
\end{tabular}
\end{center}

这在长上下文推理时尤为关键：
处理 128K token 的上下文时，
MLA 的 KV cache 仅为标准 MHA 的约 3\%，
使得在单机上处理超长上下文成为可能。
DeepSeek-V3 和后续版本都采用了 MLA，
这是 DeepSeek 系列模型在推理效率上领先的关键因素之一。

MLA 的成功也为注意力机制的设计提供了一个更一般的启示：
\textbf{KV 空间中存在大量可利用的冗余}。
如果模型在 16384 维的 KV 空间中
有效信息只占 512 维的子空间，
那么存储和计算全部 16384 维就是巨大的浪费。
MLA 的低秩投影正是对这种冗余的显式利用：
它也暗示了未来可能出现更极端的压缩方案，
甚至可以根据输入内容\textbf{自适应地}选择压缩维度。

\begin{warnbox}[MLA 的限制]
MLA 的低秩压缩假设了 KV 空间中的冗余性。
对于需要极高注意力精度的任务
（如需要精确回忆长文档中特定细节），
压缩可能带来信息损失。
DeepSeek 的解决方案是对 RoPE 部分单独处理：
位置编码对应的 KV 分量不经过压缩，
直接存储在 cache 中。
这增加了少量额外存储，
但确保了位置信息的完整性。
\end{warnbox}

\section{mHC：流形约束超连接}
\label{sec:mhc}

2025 年 12 月 31 日，DeepSeek 发表了一篇
可能对下一代模型架构产生深远影响的论文：
\textbf{mHC}（Manifold-Constrained Hyper-Connections）
~\cite{mhc2025}（arXiv:2512.24880），
由 DeepSeek CEO 梁文锋联合署名。
这篇论文瞄准了深度学习中一个
被使用了近十年却鲜少被质疑的基础组件：
\textbf{残差连接}（residual connections）。

\subsection{残差连接的隐患}

残差连接自 ResNet（2015）以来一直是
深度网络的标配。
其核心思想 $x' = x + f(x)$ 保证了梯度
可以通过恒等映射 $I$ 畅通无阻地反向传播，
使得训练数百层深的网络成为可能。
当代所有主流 LLM 的每一个 Transformer 块
都依赖残差连接。

但 DeepSeek 的研究揭示了一个
此前被忽视的严重问题：
当残差连接被推广为
\textbf{超连接}（Hyper-Connections）：
即允许跨层的加权连接而非简单的 $x + f(x)$：
\textbf{信号放大效应会失控}。
在不加约束的超连接设计中，
连接权重矩阵的谱范数可以自由增长，
导致信号在深层网络中被逐层放大。

\subsection{3000x 信号放大问题}

论文中的实验数据令人触目惊心：
在一个 27B 参数的模型中，
使用\textbf{无约束超连接}后，
信号振幅在前向传播过程中
被放大了超过 \textbf{3000 倍}：
这种爆炸式增长直接摧毁了训练的稳定性，
模型完全无法收敛。

问题的数学根源在于：
超连接的权重矩阵 $W_{\mathrm{conn}}$ 是一般矩阵，
其特征值可以大于 1。
经过 $L$ 层传播后，信号的放大倍数
最坏情况下与 $\|W_{\mathrm{conn}}\|^L$ 成正比：
当 $L = 100$ 且 $\|W_{\mathrm{conn}}\| = 1.1$ 时，
放大倍数约为 $1.1^{100} \approx 13{,}781$ 倍。

\subsection{Sinkhorn-Knopp 约束}

mHC 的核心解决方案是
将超连接的权重矩阵约束在
\textbf{双随机矩阵流形}上。

双随机矩阵（doubly-stochastic matrix）的定义是：
所有元素非负，且每行之和与每列之和都等于 1。
这类矩阵的谱范数恰好等于 1：
意味着信号经过矩阵变换后
\textbf{不会被放大}，只会被重新分配。

从信息论的角度看，双随机矩阵描述了一个
\textbf{保持总信号能量不变}的通道：
信号的总功率在变换前后不变，
只是被重新分配到不同的通道上。
这与标准残差连接 $x + f(x)$ 的「加法」模型截然不同：
后者允许信号能量单调增长
（每层都在原始信号上\emph{叠加}新信息），
而双随机超连接则在保持总能量不变的前提下
\emph{重新路由}信息流。

将一般矩阵投影到双随机矩阵流形上的方法是
经典的 \textbf{Sinkhorn-Knopp 算法}：
交替对行和列做归一化：
\[
  W \leftarrow D_1 W D_2, \quad
  \text{where } D_1 = \mathrm{diag}(W\mathbf{1})^{-1}, \;
  D_2 = \mathrm{diag}(W^\top\mathbf{1})^{-1}
\]
交替执行行归一化（乘以 $D_1$）和列归一化（乘以 $D_2$），
几次迭代（通常 3--5 次）即可收敛到满足约束的矩阵。
这个投影操作在每次前向传播时执行，
计算开销极小，对于一个 $k \times k$ 的连接矩阵
（$k$ 通常为 4--8，远小于模型维度），
Sinkhorn 迭代的开销与一次矩阵乘法相比微不足道。

\subsection{实验效果}

mHC 在 3B、9B、27B 三个规模的模型上进行了验证：

\begin{itemize}[noitemsep]
  \item \textbf{信号放大控制}：从 3000x 降低到仅 \textbf{1.6x}：
        几乎完全消除了不受控的信号放大
  \item \textbf{推理精度提升}：在 BIG-Bench Hard 推理基准上
        获得约 \textbf{+7\%} 的准确率提升
  \item \textbf{训练开销极低}：mHC 的额外计算开销
        仅为总训练成本的约 \textbf{6.7\%}：
        Sinkhorn-Knopp 投影的计算量
        相比注意力和 FFN 的矩阵乘法微不足道
  \item \textbf{27B 模型成功训练}：
        在无约束超连接完全失败的 27B 规模下，
        mHC 保证了训练的稳定收敛
\end{itemize}

\begin{corebox}[mHC 的架构意义]
mHC 的重要性在于它质疑了一个看似“已解决”的基础问题。
残差连接从 2015 年沿用至今几乎未被修改，
而 mHC 证明了更灵活的跨层连接模式
（在正确约束下）可以显著提升模型能力。
考虑到 DeepSeek CEO 梁文锋亲自参与了这篇论文，
mHC 很可能是\textbf{未来 DeepSeek V4 架构的核心组件之一}：
用流形约束的超连接替代简单的残差连接，
让深层网络中的信息流动更加高效。
注意：截至 2026 年 3 月，DeepSeek V4 尚未正式发布，
mHC 在 V4 中的具体应用方式仍有待确认。
\end{corebox}

\section{完整的 Transformer 块}

让我们把所有组件组装在一起，
看一个完整的现代 Transformer 块（如 LLaMA 3 风格）是什么样的：

\begin{lstlisting}[caption={A complete modern Transformer block}]
class TransformerBlock(nn.Module):
    def __init__(self, d_model, n_heads, n_kv_heads, ffn_dim):
        super().__init__()
        self.attn_norm = RMSNorm(d_model)
        self.attention = GQAttention(d_model, n_heads, n_kv_heads)
        self.ffn_norm = RMSNorm(d_model)
        self.ffn = SwiGLUFFN(d_model, ffn_dim)

    def forward(self, x, freqs_rope, mask=None):
        # Pre-Norm + GQA + Residual
        h = x + self.attention(
            self.attn_norm(x), freqs_rope, mask
        )
        # Pre-Norm + SwiGLU FFN + Residual
        out = h + self.ffn(self.ffn_norm(h))
        return out

class LLM(nn.Module):
    def __init__(self, vocab_size, d_model, n_layers, ...):
        super().__init__()
        self.embed = nn.Embedding(vocab_size, d_model)
        self.layers = nn.ModuleList([
            TransformerBlock(d_model, ...) for _ in range(n_layers)
        ])
        self.norm = RMSNorm(d_model)
        self.head = nn.Linear(d_model, vocab_size, bias=False)

    def forward(self, tokens):
        x = self.embed(tokens)
        freqs = precompute_rope_freqs(...)
        for layer in self.layers:
            x = layer(x, freqs)
        x = self.norm(x)
        logits = self.head(x)
        return logits
\end{lstlisting}

这段代码展示了一个完整 LLM 的结构：
词嵌入层将 token ID 映射为向量，
经过多个 Transformer 块（每个块包含 RMSNorm + GQA + RMSNorm + SwiGLU FFN），
最后通过线性层输出词表维度的 logits。
这个 logits 经过 softmax 就是模型预测的下一个 token 的概率分布。

值得注意的是代码中的几个细节：
\begin{itemize}[noitemsep]
  \item \texttt{freqs\_rope} 是预计算的 RoPE 频率张量，
        在所有层之间共享——RoPE 的频率只取决于位置，
        不需要为每层重新计算
  \item 最终输出前还有一个 \texttt{RMSNorm}：
        这是为了稳定最后一层的输出尺度，
        确保 logits 的数值范围合理
  \item \texttt{bias=False}：现代 LLM 几乎不使用偏置项，
        因为偏置对模型性能的贡献极小，
        但会影响张量并行的实现效率
\end{itemize}

\subsection{参数量分析}

一个典型的 7B 参数模型
（如 LLaMA 3-8B，实际 8.03B 参数）的配置为：
$d_{\mathrm{model}} = 4096$，$n_{\mathrm{layers}} = 32$，
$n_{\mathrm{heads}} = 32$，$n_{\mathrm{kv\_heads}} = 8$（GQA），
$d_{\mathrm{ffn}} = 14336$（SwiGLU 的 $8d/3$），
词表 128K。

将这些数字代入，
让我们逐项计算参数量：

\textbf{嵌入层}：$128{,}000 \times 4096 \approx 0.52B$

\textbf{每个 Transformer 块}：
\begin{itemize}[noitemsep]
  \item 注意力层 $Q$ 投影：$4096 \times 4096 = 16.8M$
  \item 注意力层 $K$ 投影（GQA，8 头）：$4096 \times (8 \times 128) = 4.2M$
  \item 注意力层 $V$ 投影（同 $K$）：$4.2M$
  \item 输出投影 $W^O$：$4096 \times 4096 = 16.8M$
  \item FFN $W_1$：$4096 \times 14336 = 58.7M$
  \item FFN $W_3$（SwiGLU gate）：$4096 \times 14336 = 58.7M$
  \item FFN $W_2$：$14336 \times 4096 = 58.7M$
  \item 2 个 RMSNorm（$\gamma$ 参数）：$2 \times 4096 = 8K$
  \item 每块合计：$\sim 218M$
\end{itemize}

\textbf{32 层合计}：$32 \times 218M \approx 6.98B$

\textbf{输出层}：LLaMA 3-8B 使用\textbf{非共享}（untied）嵌入，
输出投影层有独立的参数：$128{,}000 \times 4096 \approx 0.52B$

\textbf{总计}：$0.52B\text{（嵌入）} + 6.98B\text{（Transformer 层）} + 0.52B\text{（输出头）} + 0.01B\text{（RMSNorm 等）} \approx 8.03B$

这与 LLaMA 3-8B 公布的 8.03B 参数吻合。
注意：如果嵌入层与输出层共享权重（tied embeddings），
参数量将减少约 0.5B 至 $\sim$7.5B。
LLaMA 3 选择 untied embeddings 以获得更好的模型质量。

\subsection{三种架构风格}

虽然代码展示的是 decoder-only 架构（GPT 风格），
但 Transformer 实际上有三种主要变体：

\textbf{Encoder-only（BERT 风格）}：
只有 encoder，使用双向注意力：
每个 token 可以同时关注前面和后面的 token。
适合理解型任务（分类、命名实体识别、问答）。
通过 Masked Language Model（MLM）预训练。
代表：BERT、RoBERTa、ALBERT。

\textbf{Encoder-Decoder（T5 风格）}：
encoder 处理输入序列（双向注意力），
decoder 生成输出序列（单向注意力 + 对 encoder 的交叉注意力）。
适合序列到序列任务（翻译、摘要）。
代表：T5、BART、原始 Transformer。

\textbf{Decoder-only（GPT 风格）}：
只有 decoder，使用因果注意力（causal attention）：
每个 token 只能关注它自己和前面的 token。
通过 next-token prediction 预训练。
代表：GPT 系列、LLaMA、DeepSeek、Claude。

\begin{whybox}[为什么 decoder-only 最终胜出？]
decoder-only 架构在生成任务上胜出有几个关键原因：
\begin{enumerate}[noitemsep]
  \item \textbf{训练效率}：在 next-token prediction 中，
        序列中的每个位置都是一个训练样本：
        $n$ 个 token 提供 $n-1$ 个训练信号。
        而 MLM 只对随机遮盖的 15\% token 计算损失。
  \item \textbf{统一的接口}：所有任务都可以表述为
        「给定上文，生成下文」：
        不需要为不同任务设计不同的头（head）。
        理解和生成使用同一个模型，同一组参数。
  \item \textbf{涌现能力}：in-context learning、
        chain-of-thought reasoning 等涌现能力
        天然适合自回归生成的范式。
  \item \textbf{规模优势}：在大规模预训练中，
        decoder-only 的简单性使得工程实现更高效，
        更容易扩展到万亿参数。
\end{enumerate}
这并不意味着 encoder 架构消亡了：
BERT 风格的模型在嵌入（embedding）和检索任务上仍然活跃，
GTE、E5、BGE 等现代嵌入模型仍采用 encoder 架构。
但对于通用语言模型，decoder-only 是 2024--2026 年的唯一选择。

值得一提的是，这种统治地位并非理论上的必然。
decoder-only 的胜出很大程度上是\textbf{工程规模化效应}：
当 GPT-3 首先证明了 decoder-only 在极大规模下的涌现能力后，
几乎所有后续的工程投入（框架优化、推理加速、RLHF 流程）
都围绕 decoder-only 展开。
encoder-decoder 架构并未获得同等规模的工程资源验证，
因此其在超大规模下的潜力仍然是未知的。
\end{whybox}

\subsection{计算量分析：注意力 vs FFN}

在每个 Transformer 块中，注意力层和 FFN 层的计算量分布
取决于序列长度 $n$ 和模型维度 $d$。

注意力层的计算量主要来自两部分：
\begin{itemize}[noitemsep]
  \item QKV 投影：$3 \times 2nd^2 = 6nd^2$ FLOPs
        （GQA 下 KV 投影更小，但 Q 投影不变）
  \item 注意力矩阵计算：$QK^\top$ 为 $2n^2 d_k$，
        $AV$ 为 $2n^2 d_v$，合计约 $4n^2 d/h$
  \item 输出投影：$2nd^2$ FLOPs
\end{itemize}

FFN 层的计算量（SwiGLU 三个矩阵）：
$3 \times 2n \cdot d \cdot d_{\mathrm{ffn}} \approx 3 \times 2n \times d \times 2.67d \approx 16nd^2$

合计每层约 $24nd^2 + 4n^2 d/h$ FLOPs。
当 $n \ll 6dh$（对于 $d = 4096, h = 32$，即 $n \ll 768K$），
FFN 的计算量\textbf{远超注意力}：
这就是为什么在典型的训练和推理场景中，
FFN 而非注意力才是计算瓶颈。
只有在超长上下文（$n > 100K$）场景下，
注意力的 $O(n^2)$ 项才开始主导。

这个分析解释了一个常见的误解：
许多人认为注意力是 Transformer 的计算瓶颈，
但在当前主流的训练序列长度（4K--32K）下，
\textbf{FFN 消耗的计算量约是注意力的 2--3 倍}。
Flash Attention 之所以重要，
不是因为它减少了计算量（它没有），
而是因为它减少了内存传输和内存占用。

\begin{keybox}[Transformer 层的完整计算流程]
一个现代 Transformer 层（如 LLaMA 3）的前向传播：
\begin{enumerate}[noitemsep]
  \item $x_1 = x + \mathrm{GQA}(\mathrm{RMSNorm}(x))$
        \quad {\small // Pre-Norm + 分组查询注意力 + 残差连接}
  \item $x_2 = x_1 + \mathrm{SwiGLU}(\mathrm{RMSNorm}(x_1))$
        \quad {\small // Pre-Norm + SwiGLU FFN + 残差连接}
\end{enumerate}
将这样的层堆叠 32--128 次，就构成了一个完整的 LLM。
\end{keybox}

\section{Transformer 成功的根本原因}

回顾完所有组件后，让我们思考一个更深层的问题：
\textbf{为什么 Transformer 如此成功？}
在它之前，RNN（循环神经网络）和 CNN（卷积神经网络）
也可以处理序列数据，为什么被 Transformer 完全取代了？

回答这个问题不能只看架构本身：
还必须理解 Transformer 与\textbf{硬件}的深度适配。
Transformer 的核心操作是矩阵乘法（$QK^\top$、$AV$、FFN 层），
而 GPU 的 tensor core 正是为矩阵乘法而优化的。
RNN 的核心操作是顺序递归，
这在 GPU 上几乎无法高效执行。
从某种意义上说，Transformer 的成功不仅是算法的胜利，
更是「算法--硬件协同设计」的胜利。

\subsection{RNN/LSTM 的根本瓶颈}

\textbf{RNN 的瓶颈}在于\textbf{顺序性}。
RNN 逐步处理序列中的每个位置，
第 $t$ 步的计算必须等待第 $t-1$ 步完成。
这不仅使得并行化困难（训练慢），
更关键的是信息必须通过一个固定维度的\textbf{隐状态瓶颈}
在时间步之间传递：
对于长序列，早期信息会被后来的信息覆盖。

虽然 LSTM 和 GRU 通过门控机制缓解了梯度消失问题，
但信息衰减仍然是根本性的。
LSTM 的记忆容量受限于隐状态维度（通常 512--2048），
无论序列多长，所有信息都必须被压缩到这个固定大小的向量中。
这就好比要求你用一张小纸条传递整本书的内容：
必然会丢失大量信息。

更严重的是训练速度问题。
在 RNN 中，$n$ 个时间步的计算是\textbf{串行}的，
无法被 GPU 的数千个核心并行处理。
对于一个 100K token 的序列，
RNN 需要 100K 步串行计算，
而 Transformer 的注意力可以在一步中完成所有位置的交互。
这个并行化优势使得 Transformer 在相同硬件上的训练速度
比 RNN 快\textbf{一到两个数量级}。

一个被低估的事实是：Transformer 的硬件利用率远高于 RNN。
现代 GPU（如 A100）有 312 TFLOPS 的 FP16 峰值算力，
但这些算力只有在执行\textbf{大矩阵乘法}时才能被充分利用。
Transformer 的核心操作恰好都是大矩阵乘法：
$QK^\top$ 是 $(n, d) \times (d, n)$，FFN 是 $(n, d) \times (d, 4d)$。
而 RNN 的核心操作是 $(d, d)$ 的隐状态更新：
矩阵太小，GPU 的 tensor core 大部分时间在空转等待数据。

\subsection{CNN 用于序列建模的局限}

\textbf{CNN 的瓶颈}在于\textbf{局部性}。
卷积操作天然具有固定的感受野：
每个位置只能「看到」前后若干个位置。
要捕捉远距离依赖，需要堆叠很多层
（每层扩大感受野），效率低下。

一个卷积核大小为 $k$ 的 CNN，
需要 $O(n/k)$ 层才能让两个相距 $n$ 的位置直接交互。
对于 $n = 128K$、$k = 5$ 的情况，需要约 25K 层：
这在实践中完全不可行。

虽然有扩张卷积（dilated convolution）等技巧可以加速感受野增长，
但 CNN 始终缺乏在两个任意位置之间建立直接连接的能力。

\subsection{Transformer 的优势总结}

\textbf{Transformer 的优势}在于它直接建模
\textbf{任意两个位置之间的关系}：
注意力机制允许位置 1 直接与位置 1000 交互，
不需要中间位置作为「中继」。
而且注意力的计算是完全\textbf{并行}的：
所有位置同时参与，充分利用 GPU 的并行计算能力。

\begin{center}\small
\renewcommand{\arraystretch}{1.3}
\begin{tabular}{lccc}
\toprule
\rowcolor{figLightGray}
\textbf{特性} & \textbf{RNN/LSTM} & \textbf{CNN} & \textbf{Transformer} \\
\midrule
最大路径长度 & $O(n)$ & $O(\log_k n)$ & $O(1)$ \\
训练并行度 & 低（串行） & 高（并行） & 高（并行） \\
长距离依赖 & 差（信息衰减） & 中（需要深网络） & 强（直接注意力） \\
每层计算量 & $O(n \cdot d^2)$ & $O(k \cdot n \cdot d^2)$ & $O(n^2 \cdot d)$ \\
\bottomrule
\end{tabular}
\end{center}

当然，Transformer 也有代价：$O(n^2)$ 的复杂度。
但 Flash Attention 等技术在很大程度上缓解了这个问题。
而且在实践中，Transformer 的训练效率
（单位计算量学到的知识）远高于 RNN 和 CNN，
因此即使单步计算量更大，总体效率仍然更优。

\subsection{规模化特性：为什么 Transformer 越大越好}

Transformer 成功的另一个关键原因是它\textbf{优异的规模化特性}。
Kaplan 等人~\cite{kaplan2020scaling} 发现了
Transformer 的\textbf{缩放定律}（scaling laws）：
模型的损失值随参数量、数据量和计算量的增加
呈幂律下降，而且这种下降在多个数量级上都保持稳定。

这种可预测的规模化行为在 RNN 和 CNN 中从未被观察到。
RNN 在超过一定规模后，训练不稳定性急剧上升，
梯度剪裁和学习率调整的难度使得进一步扩展变得不可靠。
CNN 用于序列建模时，感受野的对数增长限制了
模型从更深网络中获益的能力。

Transformer 的规模化优势有两个结构性原因：
\begin{enumerate}[noitemsep]
  \item \textbf{残差连接 + Pre-Norm} 保证了梯度在任意深度都能有效传播，
        使得增加层数几乎不增加训练难度
  \item \textbf{注意力的表达能力随头数和模型维度平滑增长}：
        更多的头可以学习更多类型的依赖关系，
        更大的维度提供了更丰富的表示空间
\end{enumerate}

正是这种「可靠地越大越好」的特性，
使得研究实验室和公司愿意投入巨额计算资源来训练更大的 Transformer：
因为回报是可预测的。
如果 Transformer 的性能在某个规模后饱和或波动，
LLM 的军备竞赛根本不会发生。

\subsection{Attention is All You Need 的深远影响}

Vaswani 等人~\cite{vaswani2017attention}
发表的这篇论文影响了几乎所有的 AI 领域，
远超其最初的机器翻译应用：

\begin{itemize}[noitemsep]
  \item \textbf{NLP}：BERT、GPT 系列、T5 等所有现代语言模型
  \item \textbf{计算机视觉}：Vision Transformer（ViT）证明 Transformer
        可以替代 CNN 处理图像
  \item \textbf{语音}：Whisper 等模型用 Transformer 处理音频
  \item \textbf{蛋白质结构}：AlphaFold 2 的核心是 Transformer
  \item \textbf{多模态}：GPT-4V、Gemini 等多模态模型统一处理文本、图像、音频
  \item \textbf{科学计算}：天气预报（GraphCast）、
        分子设计、代码生成等领域
\end{itemize}

这种跨领域的统治力是前所未有的：
没有哪个单一架构曾经在如此多的领域中同时成为最优选择。

这种统治力的根源在于 Transformer 的\textbf{归纳偏置的最小化}。
CNN 内置了平移不变性和局部性假设：
对图像很好，对序列数据则不然。
RNN 内置了顺序处理假设：
对某些时序任务有优势，但限制了并行化。
Transformer 几乎不内置任何领域特定假设：
注意力机制允许\emph{任意两个位置}之间的交互，
位置编码是可替换的，架构本身对输入的维度和类型无感知。
这种「最小假设」的设计哲学，
使得同一架构可以无缝迁移到文本、图像、音频、蛋白质等
性质截然不同的数据模态。

\subsection{挑战者：状态空间模型}
\label{sec:ssm-challengers}

2024--2026 年，状态空间模型（State Space Model, SSM）
从学术探索走向了工程实践，
成为 Transformer 最有力的挑战者。
以 Mamba~\cite{gu2024mamba}、RWKV~\cite{peng2025rwkv7} 为代表的
SSM 架构在长序列处理上具有 $O(n)$ 的复杂度优势，
吸引了广泛的关注和跟进工作。

\subsubsection{SSM 的核心原理}

SSM 的核心是一个连续时间的线性动力系统：
\begin{equation}
  h'(t) = A h(t) + B x(t), \quad y(t) = C h(t)
\end{equation}
离散化后变成递归形式：
$h_t = \bar{A} h_{t-1} + \bar{B} x_t$，$y_t = C h_t$。
关键在于：隐状态 $h_t$ 的维度是\textbf{固定的}，
不随序列长度增长。
处理 1K 和 100K token 使用相同大小的隐状态：
这就是 $O(n)$ 复杂度的来源。

Mamba（2023）的关键创新是\textbf{选择性机制}：
矩阵 $\bar{B}$ 和 $C$ 不再是固定参数，
而是依赖于当前输入 $x_t$ 动态生成的。
这使得模型可以选择性地记住或遗忘信息，
类似于 LSTM 的门控机制，
但具有更高的并行化潜力。

为什么选择性如此关键？
考虑一个固定的 $\bar{A}$ 矩阵（非选择性 SSM）：
它以完全相同的方式处理每一个输入 token，
无论该 token 是重要的关键词还是无意义的填充词。
这就像一个没有注意力选择能力的记忆系统：
要么记住所有东西（包括噪声），
要么按固定的衰减率遗忘所有东西。
选择性机制让 $\bar{B}$ 和 $C$ 随输入变化，
使得模型可以对重要 token「打开记忆通道」，
对无关 token「关闭记忆通道」：
这在功能上近似于注意力的选择性聚焦能力。

\subsubsection{Mamba-2：结构化状态空间对偶}

Mamba-2~\cite{dao2024mamba2}（2024）提出了
\textbf{结构化状态空间对偶}（Structured State Space Duality, SSD）理论，
揭示了 SSM 和注意力之间的深层数学联系：
两者本质上都是\textbf{结构化矩阵}的不同参数化形式。
具体而言，SSD 证明了选择性 SSM 等价于一种特殊的半可分矩阵
（semiseparable matrix），
而标准注意力等价于一种因果掩码矩阵。

这个理论联系带来了直接的工程收益：
\begin{itemize}[noitemsep]
  \item \textbf{训练速度提升 2--8$\times$}：
        Mamba-2 利用 SSD 的矩阵结构，
        将计算分解为可高效并行的块矩阵乘法，
        充分利用 GPU tensor core 的硬件特性
  \item \textbf{状态空间扩大 8$\times$}：
        Mamba-1 的状态维度为 $N = 16$，
        Mamba-2 扩大到 $N = 128$：
        更大的状态意味着更强的信息存储能力，
        部分缓解了「信息瓶颈」问题
  \item \textbf{更优的并行化}：
        Mamba-2 的 SSD 算法天然支持 tensor parallelism，
        只需每层 1 次 all-reduce（与 Transformer 相同），
        而 Mamba-1 需要 2 次
\end{itemize}

\subsubsection{RWKV-7：递归的另一条路线}

RWKV~\cite{peng2025rwkv7} 是 Linux Foundation 支持的开源项目，
走了与 Mamba 不同的技术路线。
RWKV-7（代号 “Goose”，2025 年 3 月发布）
引入了\textbf{广义 delta 规则}：
一种具有向量门控和自适应学习率的状态更新机制。

RWKV-7 的关键特性：
\begin{itemize}[noitemsep]
  \item \textbf{恒定内存和推理时间}：
        与 Mamba 相同，每个 token 的推理复杂度为 $O(1)$
        （不随序列长度增长）
  \item \textbf{状态追踪能力}：
        RWKV-7 能执行状态追踪（state tracking）
        并识别所有正则语言：
        这超越了 Transformer 在标准复杂度猜想下的 $\mathrm{TC}^0$ 限制
  \item \textbf{多语言性能}：
        2.9B 参数的 RWKV-7 在多语言任务上达到 3B 级别 SoTA，
        尽管训练 token 数量（3.1T）远少于同级别的 Transformer 模型
  \item \textbf{训练可并行化}：
        尽管是递归架构，
        RWKV-7 的训练仍然可以通过分块并行化高效进行
\end{itemize}

\subsubsection{混合架构：两者兼得}

纯 SSM 模型在需要精确回忆（exact recall）的任务上
仍然弱于 Transformer：
例如「在长文档中找到一个特定句子并逐字复制」。
这是因为固定维度的隐状态
无法像注意力的 $n \times n$ 矩阵那样
保留对每个 token 的精确访问。

\textbf{混合架构}（Hybrid Architecture）
因此成为了更实际的方向：
在一个模型中交替使用注意力层和 SSM/线性递归层：

\textbf{Jamba 1.5}（AI21 Labs，2024，ICLR 2025）~\cite{jamba2024}：
Transformer-Mamba-MoE 三者结合的混合架构。
Jamba 1.5-Large 有 94B 活跃参数（MoE 总参数更大），
Jamba 1.5-Mini 有 12B 活跃参数。
两个版本都支持 256K token 的有效上下文长度：
在发布时是开源权重模型中最长的。
Jamba 的关键设计选择是\textbf{不均匀混合}：
大约 1:7 的比例交替 Transformer 层和 Mamba 层，
Transformer 层负责全局信息整合，
Mamba 层负责高效的局部处理。
实测吞吐量在长上下文场景下显著优于纯 Transformer 模型。

\textbf{Zebra-LLaMA}（AMD，2025）~\cite{zebra2025}：
将 SSM 层和 MLA（多头潜在注意力）层交替组合，
从已有的预训练 Transformer 权重初始化，
通过少量后训练（post-training）达到接近 Transformer 的精度，
同时获得接近 SSM 的推理效率。
这种「从现有 Transformer 转换」的思路，
大幅降低了混合架构的训练成本。

\textbf{OLMo Hybrid}（AI2，2026）~\cite{olmo2026hybrid}：
用 Gated DeltaNet 层替换标准 Transformer 中的
滑动窗口注意力层，保留部分全局注意力层。
7B 参数的 OLMo Hybrid 在标准预训练和中期训练评估中
全面优于同规模的纯 Transformer 模型 OLMo 3。
更重要的是，AI2 从理论上证明了
混合架构的表达能力\textbf{严格超过}
纯 Transformer 和纯线性 RNN：
混合模型可以表达两者都无法单独表达的任务。

\begin{center}\small
\renewcommand{\arraystretch}{1.3}
\begin{tabular}{lcccc}
\toprule
\rowcolor{figLightGray}
\textbf{架构} & \textbf{推理复杂度} & \textbf{训练并行} & \textbf{精确回忆} & \textbf{代表模型} \\
\midrule
纯 Transformer & $O(n^2)$ & 高 & 强 & GPT-4o, LLaMA 3 \\
纯 SSM & $O(n)$ & 中--高 & 弱 & Mamba-2, RWKV-7 \\
混合架构 & $O(n)$--$O(n^2)$ & 高 & 中--强 & Jamba 1.5, OLMo Hybrid \\
\bottomrule
\end{tabular}
\end{center}

\subsubsection{SSM 在 frontier 模型中的采用现状}

截至 2026 年 3 月，所有 frontier 模型
（GPT-4.5、Claude、Gemini、DeepSeek-V3、LLaMA 4）
仍然使用纯 Transformer 或 Transformer-MoE 架构作为主体。
\textbf{没有任何 frontier 模型完全采用 SSM 层。}

但混合架构正在快速进入中小规模模型：
Jamba 1.5、OLMo Hybrid、Nvidia Nemotron 3 Nano、
IBM Granite 4、Qwen 3.5 等模型
都在不同程度上引入了线性递归层。
Nathan Lambert 在 2026 年 3 月的分析中指出：
混合架构正经历「到处同时被采用」的时刻~\cite{lambert2026hybrid}。

这场竞争的本质不是「SSM 替代 Transformer」，
而是\textbf{最优的注意力与递归的混合比例是什么？}
一些初步的经验规则开始浮现：
\begin{itemize}[noitemsep]
  \item 大约 10--20\% 的层使用全局注意力
        足以保持精确回忆能力
  \item SSM/递归层在处理局部依赖
        和长距离的「渐进式」信息整合时效率更高
  \item 混合架构在长上下文推理中的 KV cache 显著更小
        （SSM 层不需要 KV cache）
\end{itemize}

\subsubsection{Mamba-3：推理优先的 SSM}

2026 年初，Tri Dao 和 Albert Gu 等人发表了 Mamba-3，
获选 ICLR 2026 Oral 论文。
Mamba-3 提出了“推理优先”（inference-first）的 SSM 设计哲学，
围绕三项关键改进重新设计了 SSM 的核心计算：

\begin{itemize}[noitemsep]
  \item \textbf{改进的离散化}（improved discretization）：
        更精确地将连续时间动力系统映射到离散递归形式，
        减少了离散化引入的近似误差
  \item \textbf{复数动力学}（complex dynamics）：
        在复数域上定义状态转移矩阵 $\bar{A}$，
        使模型能够捕捉振荡模式和周期性依赖：
        这在实数域中难以高效表达
  \item \textbf{MIMO 更新}（Multiple-Input Multiple-Output）：
        允许多个输入通道同时更新多个输出通道，
        突破了 Mamba-1/2 中 SISO 结构的表达瓶颈
\end{itemize}

Mamba-3 保持了 SSM 的核心优势：
\textbf{线性计算复杂度和恒定内存占用}，
同时在语言建模质量上进一步缩小了与 Transformer 的差距。
其“推理优先”的设计理念意味着
架构选择从一开始就面向高效推理场景优化，
而非先追求训练性能再做推理适配。

\subsubsection{NVIDIA Nemotron Nano 2：工业级混合架构}

NVIDIA 在 Nemotron Nano 2 中将混合架构
从学术探索推进到了工业部署阶段。
Nemotron Nano 2 基于 \textbf{Nemotron-H 架构}：
一种 Mamba-Transformer 混合设计，
将模型中\textbf{大部分自注意力层替换为 Mamba-2 层}，
仅保留少量全局注意力层用于精确信息检索。

在 9B 参数规模下，Nemotron Nano 2 在多项标准基准上
达到了 SOTA 水平，同时推理吞吐量显著高于
同规模的纯 Transformer 模型。
这一结果表明，混合架构不仅在学术指标上有竞争力，
在真实部署场景的性能--效率权衡中
同样是更优的选择。

\subsubsection{RWKV-X：稀疏注意力与递归记忆的融合}

RWKV 系列的最新迭代 RWKV-X 探索了一种新的融合方向：
将\textbf{稀疏注意力}与\textbf{递归记忆}相结合，
实现百万 token 级别的线性复杂度解码。
不同于 Jamba 等混合架构
在层级别交替使用注意力和 SSM，
RWKV-X 在单层内部融合了两种机制：
稀疏注意力用于选择性地回忆关键信息，
递归记忆用于维护连续的上下文状态。
这种设计在 1M token 的长序列解码中
保持了线性的计算和内存复杂度。

\begin{whybox}[为什么 SSM 至今未能取代 Transformer？]
SSM 的 $O(n)$ 复杂度优势是真实的，
但面临几个实际障碍：
\begin{enumerate}[noitemsep]
  \item \textbf{硬件生态}：
        7 年来，整个 GPU 软件栈
        （CUDA 库、编译器、推理引擎）
        都围绕矩阵乘法和注意力优化。
        SSM 的扫描操作（scan）缺乏同级别的硬件支持
  \item \textbf{工程惯性}：
        训练框架（Megatron、DeepSpeed）、
        推理引擎（vLLM、TensorRT-LLM）
        都深度优化了 Transformer 的执行路径。
        支持新架构需要大量工程投入
  \item \textbf{精确回忆的差距}：
        在需要从长上下文中精确检索信息的任务上，
        注意力的 $n \times n$ 精确匹配
        仍然优于 SSM 的压缩隐状态
  \item \textbf{规模验证不足}：
        截至 2026 年初，最大的纯 SSM 模型仅有 $\sim$3B 参数，
        而 frontier Transformer 已到万亿参数规模。
        SSM 在超大规模下的表现仍是未知数
\end{enumerate}
未来的赢家大概率是混合架构：
保留注意力的精确性，
同时利用 SSM 的效率优势。
\end{whybox}

\subsection{新型注意力变体}
\label{sec:new-attention-variants}

2024--2025 年，研究者们不仅在探索 Transformer 的替代品，
也在改进注意力机制本身。
两个值得关注的方向是
\textbf{差分注意力}（Differential Attention）
和\textbf{原生稀疏注意力}（Native Sparse Attention）。

\subsubsection{差分 Transformer（DIFF Transformer）}

Microsoft Research 提出的 DIFF Transformer~\cite{ye2024diff}
（ICLR 2025 Oral）瞄准了注意力机制的一个根本缺陷：
\textbf{标准注意力会将过多的权重分配给无关 token}。

即使在 softmax 归一化后，
大量无关 token 的注意力分数虽然单个很小，
但加起来仍然构成显著的噪声。
这种「注意力噪声」导致模型容易被无关上下文干扰，
是幻觉（hallucination）和 in-context learning
对示例顺序敏感等问题的根源之一。

DIFF Transformer 的解决方案在概念上优雅而简洁：
\textbf{用两个独立的 softmax 注意力图之差来计算注意力}：
\begin{equation}
  \mathrm{DiffAttn}(Q, K, V) =
  \left(\mathrm{softmax}\!\left(\frac{Q_1 K_1^\top}{\sqrt{d}}\right)
  - \lambda \cdot \mathrm{softmax}\!\left(\frac{Q_2 K_2^\top}{\sqrt{d}}\right)\right) V
\end{equation}
其中 $Q_1, Q_2, K_1, K_2$ 分别是查询和键的两个独立投影，
$\lambda$ 是可学习的标量参数。

这种差分操作的物理直觉类似于\textbf{降噪}：
类比于噪声消除耳机的工作原理：
一个麦克风录入环境噪声，
另一个录入信号+噪声，
两者相减就消除了共同的噪声成分。
在 DIFF Transformer 中，
两个注意力头分别捕捉「信号+噪声」和「噪声」，
相减后保留下来的是真正有区分度的注意力信号。
结果是注意力分布变得更加稀疏和精确，
模型更能聚焦于真正相关的 token。

$\lambda$ 参数的设计也值得注意：
它不是一个全局常数，而是\textbf{每层可学习}的。
在实践中，浅层的 $\lambda$ 通常较小
（差分程度弱，保留更多信息），
深层的 $\lambda$ 较大
（差分程度强，更精确地聚焦关键信息）。
这种层间变化暗示了注意力在不同深度
扮演着不同的角色：
浅层做广泛的信息收集，深层做精确的信息筛选。

实验表明 DIFF Transformer 的优势：
\begin{itemize}[noitemsep]
  \item \textbf{参数效率}：
        DIFF Transformer 仅需约 65\% 的模型参数
        或训练 token 就能达到与标准 Transformer 相当的语言建模性能
  \item \textbf{幻觉缓解}：
        在问答和摘要任务中显著降低幻觉率：
        因为模型更少被无关上下文「误导」
  \item \textbf{ICL 鲁棒性}：
        in-context learning 的准确率更高，
        且对示例排列顺序的敏感度大幅降低：
        这曾被认为是注意力机制的顽疾
  \item \textbf{长上下文}：
        在长上下文建模和关键信息检索任务上表现更优
  \item \textbf{激活异常值减少}：
        注意力输出中的异常值更少，有利于量化和高效推理
\end{itemize}

后续工作 DEX~\cite{kong2025dex}（NeurIPS 2025）
进一步证明了差分注意力的优势可以
\textbf{廉价地移植到已有的预训练模型}：
通过在输出值矩阵上添加轻量级差分操作，
仅需不到 0.01\% 的适应数据
即可显著提升预训练 LLM 的性能。

\subsubsection{原生稀疏注意力（NSA）}

DeepSeek 提出的
原生稀疏注意力（Native Sparse Attention, NSA）~\cite{yuan2025nsa}
（ACL 2025）从另一个角度优化注意力：
\textbf{不是修改注意力的计算方式，而是减少参与注意力计算的 token 数量}。

传统的稀疏注意力方法（如 Longformer、BigBird）
通常在推理时通过固定模式（如滑动窗口 + 全局 token）
跳过部分注意力计算。
但这些方法有两个根本问题：
（1）固定的稀疏模式无法适应不同的输入内容；
（2）训练时仍然使用全注意力，
训练和推理的不一致导致质量损失。

NSA 的核心创新是\textbf{动态分层稀疏策略}，
结合了三种互补的注意力分支：
\begin{enumerate}[noitemsep]
  \item \textbf{压缩分支}（Compression）：
        将连续的 token 块压缩为粗粒度的摘要表示，
        保持全局上下文感知
  \item \textbf{选择分支}（Selection）：
        基于重要性评分动态选择关键 token，
        实现细粒度的信息检索
  \item \textbf{滑动窗口分支}（Sliding Window）：
        保留局部上下文的完整注意力，
        确保短距离依赖不受影响
\end{enumerate}

关键的工程突破是 NSA 可以\textbf{端到端训练}：
不需要先用全注意力训练再蒸馏到稀疏模型。
在 64K 长度的序列上，NSA 在解码、前向传播和反向传播
三个阶段都实现了对全注意力的显著加速，
同时在通用基准、长上下文任务和指令推理上
保持或超过全注意力模型的性能。

\begin{corebox}[注意力机制的演化谱系]
从 2017 年的标准注意力到 2025 年，
注意力机制已经形成了丰富的变体谱系：
\begin{itemize}[noitemsep]
  \item \textbf{MHA → GQA → MLA}：减少 KV cache 内存
  \item \textbf{Flash Attention 1 → 2 → 3}：利用 GPU 内存层次加速
  \item \textbf{DIFF Transformer}：差分降噪，提升注意力精度
  \item \textbf{NSA}：动态稀疏，减少参与计算的 token 数量
  \item \textbf{iRoPE}：交替有/无位置编码，改善长上下文外推
\end{itemize}
这些创新不是互斥的：
NSA 可以与 MLA 结合（如 DeepSeek-V3.2 的 MLA + DSA），
DIFF Transformer 可以与 Flash Attention 集成。
未来的模型很可能会组合多种注意力优化。
\end{corebox}

\subsubsection{Nexus：高阶注意力}

Huawei 在 2025 年 12 月提出的
Nexus~\cite{nexus2025}（arXiv:2512.03377）
从数学基础层面重新审视了注意力机制的表达瓶颈。
标准 softmax 注意力的核心操作是 $QK^\top$：
一个\textbf{双线性}运算，其输出矩阵的秩
受限于 $\min(d_k, n)$。
当 $d_k$ 远小于 $n$（在长序列中常见），
注意力矩阵是\textbf{低秩}的，
无法表达复杂的多 token 交互模式。

Nexus 通过\textbf{递归框架}实现高阶注意力：
将标准的二阶交互（$Q$ 与 $K$ 的点积）
扩展为可配置阶数的高阶交互。
直觉上，二阶注意力回答的是
“token A 与 token B 有多相关？”，
而高阶注意力可以同时建模
“token A、B、C 三者之间的关系是什么？”。
这种多元交互能力使得
注意力矩阵突破了低秩瓶颈，
在需要复杂推理的任务上展现出显著优势。

\subsubsection{Google Sequential Attention}

Google 在 2026 年初提出了 \textbf{Sequential Attention}：
一种将注意力计算分解为\textbf{有序阶段}的新范式。
不同于标准注意力一次性计算所有 token 之间的交互，
Sequential Attention 将信息处理分为多个逐步深化的阶段：
每个阶段处理不同粒度或不同类型的依赖关系，
后续阶段在前一阶段的输出基础上进一步精化。

这种分阶段设计带来了两个实际优势：
一是推理速度更快：
早期阶段可以快速筛选相关信息，
避免后续阶段对所有 token 做昂贵的精细计算；
二是上下文质量更高：
分层处理使模型在长序列中更好地保持信息的结构化组织，
减少了单一注意力层中的信息混淆。
Google 预计将在 2026--2027 年间
将 Sequential Attention 集成到 Gemini 和 Search 产品线中。

\section{词嵌入：让词进入向量空间}

在 Transformer 的所有组件之前，
还有一个基础但至关重要的步骤：
\textbf{将离散的 token ID 转换为连续的向量}。
这个过程由\textbf{嵌入层}（Embedding Layer）完成。

嵌入层本质上是一个查找表：
一个 $V \times d$ 的矩阵（$V$ 是词表大小，$d$ 是模型维度），
每一行是一个 token 的向量表示。
给定 token ID $t$，嵌入向量就是矩阵的第 $t$ 行。

值得注意的是，这里的「词」并非自然语言中的词。
现代 LLM 使用\textbf{子词分词}（subword tokenization）：
常见词保留完整形式，
罕见词被拆分为更小的子词单元。
例如，BPE（Byte Pair Encoding）算法
通过统计训练语料中的字符对频率，
迭代地合并最频繁的字符对来构建词表。
LLaMA 3 使用 128K 大小的 BPE 词表，
这个大小的选择是效率与表达能力的权衡：
词表越大，每个文本需要的 token 数越少（推理更快），
但嵌入层的参数量也越大。

这些嵌入向量在训练过程中通过梯度下降学习。
训练结束后，语义相近的词在向量空间中会距离较近：
「国王」和「女王」的嵌入向量之间的距离
小于「国王」和「苹果」之间的距离。
更有趣的是，这些向量之间会出现有意义的\textbf{方向}：
经典的例子是 $\vec{\text{King}} - \vec{\text{Man}} + \vec{\text{Woman}} \approx \vec{\text{Queen}}$，
说明嵌入空间捕获了「性别」的概念方向。

在一些 LLM 中，嵌入层和最后的输出层
（将隐状态映射回词表维度）\textbf{共享权重}：
这被称为 tied embeddings。
直觉上，一个 token 的「输入表示」和「输出表示」
应该是一致的。权重共享可以减少参数量
（对于 128K 词表、4096 维模型，节省约 0.5B 参数）。
但近年来的趋势是使用\textbf{非共享}（untied）嵌入：
例如 LLaMA 3 系列就使用独立的输出投影层，
实验表明这在大词表场景下能获得更好的模型质量。

\subsection{嵌入空间的几何结构}

训练完成后的嵌入空间具有丰富的几何结构，
远不止「相似的词距离近」这么简单。

\textbf{各向异性}（anisotropy）是嵌入空间的一个重要特征。
Ethayarajh~\cite{ethayarajh2019contextual} 发现，
在 BERT 等模型中，嵌入向量倾向于占据向量空间的一个
\textbf{窄锥}（narrow cone）而非均匀分布。
这意味着任意两个 token 的嵌入向量的余弦相似度
平均值显著大于零（通常约 0.5--0.8），
而不是接近零（随机向量在高维空间中的期望）。

这种各向异性带来了一个实际问题：
如果所有向量都在一个窄锥内，
它们之间的\textbf{相对差异}会被压缩，
使得语义距离的度量变得不精确。
后续的各种校准技术（如 whitening、isotropy calibration）
正是为了解决这个问题。

更有趣的是上下文化嵌入的行为。
Transformer 输出的每个位置的向量
不再是固定的词嵌入：
同一个词在不同上下文中会有不同的表示。
例如，「bank」在「river bank」和「bank account」中
的上层隐状态向量截然不同。
这种\textbf{上下文化}是 Transformer 表达能力的核心体现，
也是为什么 Transformer 的表示能力远超静态词向量（如 Word2Vec）的原因。

\begin{whybox}[嵌入层是否需要缩放？]
在原始 Transformer 中，嵌入向量会乘以 $\sqrt{d}$。
原因是：随机初始化的嵌入向量的模长（$L^2$ 范数）与 $\sqrt{d}$ 成正比，
而位置编码的模长固定为 1。
如果不缩放，位置信息会在嵌入信息中被淹没。
乘以 $\sqrt{d}$ 使两者在同一数量级。
现代模型（使用 RoPE 而非加法位置编码）通常不需要这个缩放，
因为 RoPE 直接作用于注意力计算而非嵌入空间。
\end{whybox}

\section{因果掩码：让模型「不看未来」}

在 decoder-only 模型中，有一个至关重要但容易被忽视的组件：
\textbf{因果掩码}（causal mask）。

自回归语言模型的训练目标是预测下一个 token。
如果 token $i$ 在计算注意力时能看到 token $i+1, i+2, \ldots$ 的信息，
那预测就变成了「抄答案」，模型不需要学到任何有意义的模式。
因果掩码通过在注意力分数矩阵上施加一个上三角掩码来解决这个问题：
\[
  \mathrm{Mask}_{ij} = \begin{cases}
    0 & \text{if } j \leq i \text{（可以看到）} \\
    -\infty & \text{if } j > i \text{（不能看到）}
  \end{cases}
\]
将掩码加到注意力分数上后，$j > i$ 的分数变为 $-\infty$，
经过 softmax 后权重变为 0：
token $i$ 的注意力只分布在位置 $1$ 到 $i$ 上。

因果掩码带来了一个重要的训练优势：
在一次前向传播中，序列中的每个位置 $i$
都在预测位置 $i+1$ 的 token，
\textbf{$n$ 个 token 提供 $n-1$ 个训练信号}。
相比之下，BERT 的 MLM 预训练只对 15\% 的被遮盖 token 计算损失，
训练效率低约 6 倍。

\subsection{因果掩码与训练并行性}

因果掩码的另一个好处是它使得训练\textbf{完全可并行}。
尽管因果注意力在概念上是「从左到右」的，
但由于掩码在训练前就已确定（不依赖于模型输出），
所有位置的注意力可以在一次矩阵运算中同时计算。
$QK^\top + \mathrm{Mask}$ 的计算不需要等待任何位置的输出：
整个序列可以在一步中并行处理。

这与推理（inference）时的情况形成鲜明对比。
在自回归生成中，每个新 token 的生成依赖于之前所有 token，
因此必须逐个生成，这正是 KV cache 被发明来加速的原因。
训练时的并行性和推理时的串行性，
是理解 Transformer 工程优化的关键区分。

\subsection{前缀缓存与灵活掩码}

在现代 LLM 的部署中，因果掩码的严格限制有时会被放松。
\textbf{前缀模式}（prefix mode）允许一段前缀
使用\textbf{双向注意力}（每个 token 可以看到整个前缀），
只有生成部分使用因果掩码。
这种混合掩码策略在 chat 模型中很常见：
系统提示词（system prompt）使用双向注意力进行充分理解，
用户的回复部分则使用因果掩码逐步生成。

\begin{keybox}[本章回顾：Transformer 的核心组件]
将本章讨论的所有组件汇总，
一个现代 LLM（如 LLaMA 3）的 Transformer 由以下核心部分组成：
\begin{enumerate}[noitemsep]
  \item \textbf{词嵌入}：BPE 分词 + 嵌入查找表，将 token ID 映射为 $d$ 维向量
  \item \textbf{RoPE 位置编码}：通过分块旋转矩阵编码相对位置
  \item \textbf{GQA 注意力}：分组共享 KV 的多头注意力，平衡质量与效率
  \item \textbf{因果掩码}：阻止模型看到未来 token
  \item \textbf{Flash Attention}：IO-aware 的分块计算，$O(n)$ 内存
  \item \textbf{SwiGLU FFN}：门控前馈网络，充当知识存储
  \item \textbf{RMSNorm}：仅做尺度归一化的 Pre-Norm
  \item \textbf{残差连接}：保证深层网络的梯度流畅
\end{enumerate}
这些组件不是独立的，它们相互配合形成一个紧密耦合的系统。
RoPE 与注意力机制耦合，Pre-Norm 与残差连接协同，
SwiGLU 的三个矩阵配合门控机制工作。
理解每个组件的「为什么」和组件之间的「怎么配合」，
是理解 LLM 架构设计的关键。
\end{keybox}

%% ============================================================
%% NEW REFERENCES (to be added to refs.bib) — includes expansion references:
%% ainslie2023gqa = Ainslie, Joshua et al. "GQA: Training Generalized Multi-Query Transformer Models from Multi-Head Checkpoints." EMNLP 2023. arXiv:2305.13245.
%% clark2019bert = Clark, Kevin et al. "What Does BERT Look At? An Analysis of BERT's Attention." ACL 2019 Workshop BlackboxNLP.
%% shazeer2019mqa = Shazeer, Noam. "Fast Transformer Decoding: One Write-Head is All You Need." arXiv:1911.02150, 2019.
%% press2022alibi = Press, Ofir et al. "Train Short, Test Long: Attention with Linear Biases Enables Input Length Generalization." ICLR 2022.
%% chen2023extending = Chen, Shouyuan et al. "Extending Context Window of Large Language Models via Positional Interpolation." arXiv:2306.15595, 2023.
%% bloc97ntk = bloc97. "NTK-Aware Scaled RoPE." Reddit/GitHub, 2023.
%% hendrycks2016gelu = Hendrycks, Dan and Gimpel, Kevin. "Gaussian Error Linear Units (GELUs)." arXiv:1606.08415, 2016.
%% ramachandran2017swish = Ramachandran, Prajit et al. "Searching for Activation Functions." arXiv:1710.05941, 2017.
%% shah2024flashattention3 = Shah, Jay et al. "FlashAttention-3: Fast and Accurate Attention with Asynchrony and Low-precision." arXiv:2407.08608, 2024.
%% gu2024mamba = Gu, Albert and Dao, Tri. "Mamba: Linear-Time Sequence Modeling with Selective State Spaces." arXiv:2312.00752, 2023 (COLM 2024).
%%
%% --- ADDED 2026-03-08 ---
%% dao2024mamba2 = Dao, Tri and Gu, Albert. "Transformers are SSMs: Generalized Models and Efficient Algorithms Through Structured State Space Duality." ICML 2024. arXiv:2405.21060.
%% peng2025rwkv7 = Peng, Bo et al. "RWKV-7 'Goose' with Expressive Dynamic State Evolution." arXiv:2503.14456, 2025.
%% jamba2024 = Jamba Team. "Jamba: A Hybrid Transformer-Mamba Language Model." ICLR 2025. arXiv:2403.19887, 2024.
%% zebra2025 = Yang, Mingyu et al. "Zebra-Llama: Towards Extremely Efficient Hybrid Models." NeurIPS 2025. OpenReview, 2025.
%% olmo2026hybrid = Merrill, William et al. "OLMo Hybrid: From Theory to Practice." Allen Institute for AI, 2026.
%% lambert2026hybrid = Lambert, Nathan. "OLMo Hybrid and Future LLM Architectures." Interconnects AI, March 2026.
%% ye2024diff = Ye, Tianzhu et al. "Differential Transformer." ICLR 2025 Oral. arXiv:2410.05258, 2024.
%% kong2025dex = Kong, Chaerin et al. "Understanding Differential Transformer Unchains Pretrained Self-Attentions." NeurIPS 2025.
%% yuan2025nsa = Yuan, Jingyang et al. "Native Sparse Attention: Hardware-Aligned and Natively Trainable Sparse Attention." ACL 2025. arXiv:2502.11089, 2025.
%%
%% --- ADDED 2026-03-12 ---
%% mhc2025 = DeepSeek-AI. "Manifold-Constrained Hyper-Connections." arXiv:2512.24880, December 2025.
%% nexus2025 = Huawei. "Nexus: Higher-Order Attention via Recursive Framework." arXiv:2512.03377, December 2025.
%%
%% --- ADDED in ch02 expansion ---
%% sanford2023representational = Sanford, Clayton et al. "Representational Strengths and Limitations of Transformers." NeurIPS 2023.
%% michel2019heads = Michel, Paul et al. "Are Sixteen Heads Really Better than One?" NeurIPS 2019.
%% katharopoulos2020linear = Katharopoulos, Angelos et al. "Transformers are RNNs: Fast Autoregressive Transformers with Linear Attention." ICML 2020.
%% jiang2023mistral = Jiang, Albert Q. et al. "Mistral 7B." arXiv:2310.06825, 2023.
%% dauphin2017glu = Dauphin, Yann N. et al. "Language Modeling with Gated Convolutional Networks." ICML 2017.
%% geva2021transformer = Geva, Mor et al. "Transformer Feed-Forward Layers Are Key-Value Memories." EMNLP 2021.
%% ba2016layernorm = Ba, Jimmy Lei et al. "Layer Normalization." arXiv:1607.06450, 2016.
%% ioffe2015batchnorm = Ioffe, Sergey and Szegedy, Christian. "Batch Normalization: Accelerating Deep Network Training." ICML 2015.
%% xiong2020prenorm = Xiong, Ruibin et al. "On Layer Normalization in the Transformer Architecture." ICML 2020.
%% wang2022deepnet = Wang, Hongyu et al. "DeepNet: Scaling Transformers to 1,000 Layers." arXiv:2203.00555, 2022.
%% kaplan2020scaling = Kaplan, Jared et al. "Scaling Laws for Neural Language Models." arXiv:2001.08361, 2020.
%% ethayarajh2019contextual = Ethayarajh, Kawin. "How Contextual are Contextualized Word Representations?" EMNLP 2019.
%% ============================================================
